\documentclass[reqno]{amsart}
\usepackage{a4wide}
\usepackage{hyperref}

\AtBeginDocument{{\noindent\small
\emph{Electronic Journal of Differential Equations},
Vol. 2026 (2026), No. 14, pp. 1--20.\newline
ISSN: 1072-6691. URL: https://ejde.math.txstate.edu, DOI: 10.58997/ejde.2026.14}
\thanks{\copyright 2026. This work is licensed under a CC BY 4.0 license.}
\vspace{8mm}}

\begin{document}
\title[\hfilneg EJDE-2026/14\hfil Fractional $(p,q)$-Laplacian equations in $\mathbb{R}^N$]
{Multiple normalized solutions for fractional $(p,q)$-Laplacian equations
in $\mathbb{R}^N$}

\author[L. Bao, L. Wang, M. Li \hfil EJDE-2026/14\hfilneg]
{Lihong Bao, Libo Wang, Min Li}

\address{Lihong Bao \newline
School of Mathematics and statistics,
Beihua University,
Jilin, Jilin, 132013, China}
\email{15246791290@163.com}

\address{Libo Wang (corresponding author) \newline
School of Mathematics and statistics,
Beihua University,
Jilin, Jilin, 132013, China}
\email{wlb\_math@163.com} 

\address{Min Li \newline
School of Mathematics and statistics,
Beihua University, 
Jilin, Jilin, 132013, China}
\email{lm\_18179579903@163.com} 

\thanks{Submitted July 26, 2025. Published February 13, 2026.}
\subjclass[2020]{35J20, 35J62, 35J92}
\keywords{Ekeland's variational principle; fractional $(p,q)$-Laplacian equation;
\hfill\break\indent normalized solutions; concentration compactness}

\begin{abstract}
In this article, we consider the multiplicity results of normalized solutions of the
fractional $(p,q)$-Laplacian equation
\begin{gather*}
(-\Delta)_p^{s} u+(-\Delta)_q^{s} u=\lambda |u|^{p-2}u+h(\varepsilon x)f(u)+g(u)
\quad\text{in }\mathbb{R}^N,\\
\int_{\mathbb {R}^{N}}|u|^p \mathrm{d}x=a^p,
\end{gather*}
where $\varepsilon, a>0$, $0<s<1<p<q$, $sq<N<\frac{spq(q-1)}{q-p}$ and
$\lambda \in \mathbb {R}$  is a unknown Lagrange multiplier,
$h$ is a continuous positive function, $f$ and $g$ are also continuous and satisfy
some growth conditions.
When $\varepsilon$ is small enough, we show that the number of normalized solutions is
at least the number of global maximum points of $h$ according to the Ekeland's
variational principle and the concentration compactness principle.
\end{abstract}

\maketitle
\numberwithin{equation}{section}
\newtheorem{theorem}{Theorem}[section]
\newtheorem{lemma}[theorem]{Lemma}
\allowdisplaybreaks

\section{Introduction}

This article focuses on the existence of multiple normalized solutions in
$X:=W^{s,p}(\mathbb {R}^N)\cap W^{s,q}(\mathbb {R}^N)$ for the
 fractional $(p,q)$-Laplacian equation
\begin{equation}\label{eqn-(1.1)}
(-\Delta)_p^{s} u+(-\Delta)_q^{s} u=\lambda |u|^{p-2}u+h(\varepsilon x)f(u)+g(u)\quad \text{in } \mathbb {R}^N
\end{equation}
under {the restriction}
\begin{equation}\label{eqn-(1.2)}
\int_{\mathbb {R}^{N}}|u|^p \mathrm{d}x=a^p,
\end{equation}
where $\varepsilon,a>0$,  $0<s<1<p<q$, $sq<N<\frac{spq(q-1)}{q-p}$ and
$\lambda \in \mathbb {R}$ is a Lagrange multiple which is unknown. The fractional
$p$-Laplace operator $(-\Delta)_p^{s}$ and the fractional
$q$-Laplace operator $(-\Delta)_q^{s}$
are defined along a smooth function (up to a normalizing constant)
$u:\mathbb{R}^N\to \mathbb{R}$ by
$$
(-\Delta)_p^{s} u(x)=2\lim_{\varepsilon\to 0}\int_{\mathbb{R}^N\backslash B_{\varepsilon}(x)}
\frac{|u(x)-u(y)|^{p-2}(u(x)-u(y))}{|x-y|^{N+sp}}\mathrm{d}y,\quad x\in\mathbb{R}^N
$$
and
$$
(-\Delta)_q^{s} u(x)=2\lim_{\varepsilon\to 0}\int_{\mathbb{R}^N\backslash
B_{\varepsilon}(x)}\frac{|u(x)-u(y)|^{q-2}(u(x)-u(y))}{|x-y|^{N+sq}}\mathrm{d}y,\quad x\in\mathbb{R}^N,$$
respectively. The continuous functions $f,g,h$ satisfies the following conditions:
\begin{itemize}
\item[(A1)]  $f$ is odd and $\lim_{t\to 0}\frac{|f(t)|}{|t|^{m_1-1}}=\alpha_1 >0$ for some
$m_1\in(p,\min\{p+\frac{sp^2}{N},m^{*}\})$, where $m^{*}= 1+\frac{Nq(p-1)}{p(N-sq)}$;

\item[(A2)] there exist constants $C_1, C_2>0$ and $m_2\in(p,\min\{q+\frac{spq}{N},m^{*}\})$
such that $|f(t)|\le C_1+C_2|t|^{m_2-1}$, $\forall t\in \mathbb{R}$;

\item[(A3)] for all $t>0$ , $qF(t)\le f(t)t$, $F(t)$ is non-decreasing, where
$F(t)=\int_0^{t}f(s)\mathrm{d}s$.

\item[(A4)] $0<h_0=\inf_{x\in \mathbb{R}^{N}}h(x)\le \max_{x\in \mathbb{R}^{N}} h(x):=h_{\rm max}$;

\item[(A5)] $h_{\infty}:=\lim_{|x|\to \infty}h(x)<h_{\rm max}$;

\item[(A6)] $h^{-1}(h_{\rm max})={e_1,e_2,\cdots,e_l}$ with $e_1=0$ and $e_j\neq e_k$ when $j\neq k$.

\item[(A7)] $g$ is odd;

\item[(A8)] there exist $q\le\beta<\min\{q+\frac{spq}{N},m^{*}\}$, such that $0\le q G(s)\le g(s)s\le\beta G(s)$, $\forall s\in\mathbb{R}$.

\item[(A9)]  $G(t)$ is non-decreasing, where $G(t)=\int_0^{t}g(s)\mathrm{d}s$.
\end{itemize}
Because of  the restriction  \eqref{eqn-(1.2)}, we are seeking a normalized solution to
\eqref{eqn-(1.1)}, which is equivalent to finding the critical point of the  functional
\begin{align*}
&I_{\varepsilon}(u) \\
&=\frac{1}{p}\iint_{\mathbb{R}^{2N}}\frac{|u(x)-u(y)|^p}{|x-y|^{N+sp}}
\mathrm{d}x\mathrm{d}y
 +\frac{1}{q}\iint_{\mathbb{R}^{2N}}\frac{|u(x)-u(y)|^q}{|x-y|^{N+sq}}
 \mathrm{d}x\mathrm{d}y
 -\int_{\mathbb{R}^N}h(\varepsilon x)F(u)\mathrm{d}x-\int_{\mathbb{R}^{N}}G(u)
 \mathrm{d}x,
\end{align*}
on the sphere
\begin{equation}\label{eqn-(1.3)}
S(a):=\{u\in X:=W^{s,p}(\mathbb {R}^N)\cap W^{s,q}(\mathbb {R}^N):|u|_p^{p}=\int_{x\in \mathbb{R}^{N}}|u|^{p}\mathrm{d}x=a^{p}\},
\end{equation}
where $|\cdot|_\tau$ denotes the usual norm on $L^{\tau}(\mathbb {R}^N)$ for
$\tau \in [1,+\infty)$.
Here the fractional Sobolev space
$W^{s,p}(\mathbb {R}^N)$ and $W^{s,q}(\mathbb {R}^N)$ are defined for any $p,q > 1$
and $s \in (0, 1)$ by
$$
W^{s,p}(\mathbb {R}^N):=\{u\in L^{p}(\mathbb {R}^N):[u]_{s,p}:=(\iint_{\mathbb{R}^{N}
\times\mathbb{R}^{N}}\frac{|u(x)-u(y)|^p}{|x-y|^{N+sp}}\mathrm{d}x\mathrm{d}y)^{1/p}<\infty\}$$
and
$$
W^{s,q}(\mathbb {R}^N):=\{u\in L^{q}(\mathbb {R}^N):[u]_{s,q}
:=(\iint_{\mathbb{R}^{N}\times\mathbb{R}^{N}}\frac{|u(x)-u(y)|^q}{|x-y|^{N+sq}}\mathrm{d}x\mathrm{d}y)^{1/q}<\infty\},
$$
with the norm
$$
\|u\|_{W^{s,p}(\mathbb {R}^N)}=(|u|_{p}^p+[u]_{s,p}^{p})^{1/p}
$$
and
$$
\|u\|_{W^{s,q}(\mathbb {R}^N)}=(|u|_{q}^q+[u]_{s,q}^{q})^{1/q}.
$$
Moreover, we set $\|u\|_X=\|u\|_{W^{s,p}(\mathbb {R}^N)}+\|u\|_{W^{s,q}(\mathbb {R}^N)}$.
It is well known that $I_{\varepsilon} \in C^{1}(X,\mathbb {R})$ and
\begin{align*}
\langle I_{\varepsilon}' (u),\varphi \rangle
&= \iint_{\mathbb{R}^{2N}}\frac{|u(x)-u(y)|^{p-2}(u(x)-u(y))(\varphi(x)-\varphi(y))}{|x-y|^{N+sp}}\mathrm{d}x\mathrm{d}y\\
&\quad +\iint_{\mathbb{R}^{2N}}\frac{|u(x)-u(y)|^{q-2}(u(x)-u(y))(\varphi(x)-\varphi(y))|}{|x-y|^{N+sq}}\mathrm{d}x\mathrm{d}y\\
&\quad -\int_{ \mathbb{R}^{N}}h(\varepsilon x)f(u)\varphi \mathrm{d}x-\int_{ \mathbb{R}^{N}}g(u)\varphi\mathrm{d}x
\end{align*}
for all $u,\varphi\in X$.


The study of problems involving fractional Laplace operators has garnered significant attention
in recent years. This growing interest stems not only from its natural extension of the classical
Laplace operators but also from its capacity to model numerous novel phenomena through its distinctive nonlinear integral structure, see \cite{ref7,ref11,ref14}.
Maya Chhetri, Petr Girg, and Elliott Hollifield \cite{ref10} systematically studied a class of
nonlinear partial differential equations involving the fractional Laplacian operator of the form
\begin{gather*}
(-\Delta)^{s}u=\lambda f(x,u)\quad  \text{in }\Omega,\\
u=0  \quad \text{in }\mathbb{R}^N\backslash\Omega,
\end{gather*}
where $f:\Omega\times\mathbb{R}\to\mathbb{R}$ is a Carath\'{e}odory function.
They established a method of subsolutions and supersolutions without monotone iteration,
proving the existence of positive weak solutions for various types of sublinear nonlinearities,
including logarithmic-type terms.
Significant research efforts have also been devoted to understanding the multiplicity of
normalized solutions for fractional Laplace operators, see \cite{ref19, ref20, ref25, ref34}.
In \cite{ref8}, employing the Mountain Pass Lemma,  the authors obtained the  existence
results of  the problem
\begin{gather*}
(-\Delta)^{s_1} u(x)+(-\Delta)^{s_2} u(x)+\lambda u(x)+V(x)u(x)=g(u(x)),\quad
x\in \mathbb{R}^{d},\\
\int_{\mathbb{R}^{d}}|u(x)|^{2}\mathrm{d}x=a,
\end{gather*}
where $0<s_1<s_2<1$, $2s_1<d<\frac{2s_1s_2}{s_2-s_1}$, $a>0$.
Function $g$  satisfies  the following conditions:
\begin{itemize}
\item[(A10)] $g$ is continuous and odd;

\item[(A11)] there exits $(\alpha, \beta)\in\mathbb{R}_{+}^{2}$ satisfying
$$
2+\frac{4s_2}{d}<\alpha<\beta<\frac{2d}{d-2s_1},
$$
such that
$$
\alpha G(s)\le g(s)s\le \beta G(s),\quad\forall s\in\mathbb{R}, \quad\text{with }
G(s)=\int_0^{s}g(t)\mathrm{d}t;
$$
\item[(A12)] let
$\tilde{G}:\mathbb{R}\to\mathbb{R}$ and
$\tilde{G}(s)=\frac{g(s)s}{2}-G(s)$.; there exists $\tilde{G}'$  and
$$
\tilde{G}'(s)s\ge\alpha\tilde{G}(s),\quad\forall s\in\mathbb{R}.
$$
\end{itemize}

In recent years, fractional $p$-Laplace operators have been extensively studied,
see \cite{ref3, ref4, ref18, ref26, ref32}. Additionally, a significant contribution was
made by Vincenzo Ambrosio and Teresa Isernia  \cite{ref1}, who analyzed the problem
\begin{gather*}
\varepsilon^{sp}(-\Delta)_p^{s} u+V(x)|u|^{p-2}u=f(x,u)\quad\text{in }\mathbb{R}^N,\\
u\in W^{s,p}(\mathbb{R}^N),\\
u(x)>0\quad  x\in\mathbb{R}^N,
\end{gather*}
where $\varepsilon>0$ is a parameter, $s\in(0,1)$, $1<p<\infty$, $N>sp$.
In the subcritical case, the authors established existence, multiplicity and concentration
results for positive solutions, with small parameters through variational methods combined with Ljusternik-Schnirelmann theory.
In case of  critical nonlinearity, they obtained existence and multiplicity of solutions of the equation
\begin{gather*}
(-\Delta)_p^{s} u+V(\varepsilon x)|u|^{p-2}u=f(u)+|u|^{p_{s}^{*}-2}u\quad\text{in }\mathbb{R}^N,\\
u\in W^{s,p}(\mathbb{R}^N),\\
u(x)>0\quad x\in\mathbb{R}^N.
\end{gather*}

Several researchers have done a great deal of research in the multiplicity of
normalized solutions of fractional Laplace operators, see \cite{ref23, ref27, ref35}.
In \cite{ref36}, the following equation was considered,
\begin{gather*}
(-\Delta)_p^{s} u+Z(kx)|u|^{p-2}u=\lambda|u|^{p-2}u+[\frac{1}{|x|^{N-q}*|u|^{q}}]|u|^{q-2}u
+\sigma|u|^{p_s^{*}-2}u,\quad x\in \mathbb{R}^{N},\\
\int_{\mathbb{R}^{N}}|u(x)|^{p}\mathrm{d}x=a^{p},
\end{gather*}
where $k>0$ is small parameter and $Z:\mathbb{R}^N\to[0,\infty)$ is a continuous function.
The authors establish the existence of multiple normalized solutions through a combination
of minimization techniques, truncation methods, variational approaches,
and Lusternik-Schnirelmann category theory. In particular, the fractional $p$-Laplace
operator, which reduces to the standard $p$-Laplacian when $s = 1$, has been extensively
studied, see \cite{ref6, ref15, ref21, n2}.
Moreover, in the special case that $s = 1$ and $p = 2$, the operator coincides with the classical Laplacian, see \cite{ref12, ref28, ref29}.

Recently, research on fractional $(p,q)$-Laplacian equations has attracted a great deal of
attention, see \cite{ref9, ref30, ref37}. Vincenzo Ambrosio and Teresa Isernia \cite{ref2} consider the equation
\begin{gather*}
(-\Delta)_p^{s}u+(-\Delta)_q^{s}u=\frac{|u|^{p_s^{*}(\alpha)-2}u}{|x|^{\alpha}}+\lambda f(x,u)\quad\text{in }\Omega,\\
u=0\quad\text{in }\mathbb{R}^N\backslash\Omega,
\end{gather*}
where $0<s<1$, $1\le q<p<\frac{N}{s}$, $0\le\alpha<sp$, and
$p_s^{*}(\alpha)=\frac{p(N-\alpha)}{N-sp}$ is Hardy-Sobolev critical exponent.
They use the concentration-compactness principle and the mountain pass lemma to show
that infinitely many solutions to the equation exist and that these solutions all tend to zero
when $\lambda$ belongs to a suitable range.

Additionally, Divya Goel, Deepak Kumar and K.Sreenadh \cite{ref16} studied the problem
\begin{gather*}
(-\Delta)_p^{s_1} u+(-\Delta)_q^{s_2} u=\lambda a(x)|u|^{\delta-2}u+ b(x) |u|^{r-2}u\quad\text{in }\Omega,\\
u=0 \quad\text{in }\mathbb{R}^N\backslash\Omega,
\end{gather*}
where $1<\delta\le q\le p<r\le p_{s_1}^{*}=\frac{np}{n-s_1p}$, $0<s_2<s_1<1$, $n>ps_1$,
and $a\in L^{\frac{r}{r-\delta}}(\Omega)$, $b\in L^{\infty}(\Omega)$ are sign changing
functions. They showed the existence and regularity of multiple weak solutions to the
convex-concave problem. To the best of our knowledge, the work concerning the
normalized solutions for the fractional $(p,q)$-Laplacian problem in $\mathbb{R}^N$ has not
been seen, see \cite{n3} for the $(p,q)$-Laplacian problem, and \cite{ref33} for the discrete fractional $(p,q)$-Laplacian problem.

Motivated by the above works, this paper aims to study the existence of multiple normalized
solutions for problem \eqref{eqn-(1.1)} and \eqref{eqn-(1.2)} by employing variational
approaches and  the concentration-compactness principle from \cite{ref1, ref13}.
However, our analysis faces two main difficulties: the first is the nonlinearity of the 
fractional $(p,q)$-Laplace operator, and the second is the lack of compactness,
since the problem  is set on $\mathbb{R}^N$.

 The main result of this paper reads as follows.

\begin{theorem}\label{theorem:1.1}
Assume {\rm (A1)--(A9)} are satisfied.
 Then there exists $\varepsilon_0>0$ such that  \eqref{eqn-(1.1)} and \eqref{eqn-(1.2)}
 has at least $l$ pairs of weak solutions $(u_j,\lambda_j)$ for $0<\varepsilon<\varepsilon_0$. Moreover, $\lambda_j<0$ and $I_{\varepsilon}(u_j)<0$ for $j=1,2,\dots,l$.
\end{theorem}

This article is organized as follows.
In Section 2, we prove the compactness theorem for the autonomous case.
In Section 3, we use the compactness theorem to study the non-autonomous case.
We prove the existence of multiple solutions in Section 4 and Section 5.

\section{Autonomous case}

Firstly, we consider the existence of the normalized solution
$(u,\lambda)\in X \times\mathbb {R}$, of the equation
\begin{equation}\label{eqn-(2.1)}
\begin{gathered}
(-\Delta)_{p}^{s} u+(-\Delta)_q^{s} u=\lambda |u|^{p-2}u+\mu f(u)+g(u),\\
\int_{\mathbb {R}^{N}}|u|^p \mathrm{d}x=a^p,
\end{gathered}
\end{equation}
where $a,\mu>0$, $\lambda\in\mathbb{R}$ and $f$ satisfies (A1)--(A3), and  $g$ satisfies
(A7)--(A9). It is well known that a critical point of the functional
$$J
_\mu(u)=\frac{1}{p}\iint_{\mathbb{R}^{2N}}\frac{|u(x)-u(y)|^p}{|x-y|^{N+sp}}
\mathrm{d}x\mathrm{d}y+\frac{1}{q}\iint_{\mathbb{R}^{2N}}\frac{|u(x)-u(y)|^q}{|x-y|^{N+sq}}\mathrm{d}x\mathrm{d}y-\mu\int_{\mathbb{R}^N}F(u)\mathrm{d}x-\int_{\mathbb{R}^{N}}G(u)\mathrm{d}x
$$
is a solution to problem \eqref{eqn-(2.1)},
which is restricted to the sphere $S(a)$.
Next, we will show that problem \eqref{eqn-(2.1)} has a normalized solution.

\begin{lemma}\label{lem:2.1}
The functional $J_\mu$ restricted to $S(a)$ is bounded from below.
\end{lemma}

\begin{proof}
From  conditions (A1) and (A2), we deduce the existence of constants
$C_1, C_2>0$ such that
$$
|F(t)|\le C_1|t|^{m_1}+C_2|t|^{m_2},\quad\forall t\in\mathbb{R}.
$$
From (A8), there exists a constant $C_3$ such that
$$
|G(t)|\le C_3(|t|^{q}+|t|^{\beta}),\quad\forall t\in\mathbb{R}.
$$
By the fractional Gagliardo-Nirenberg inequality \cite{ref24}, we have
\begin{equation}\label{eqn-(2.2)}
|u|_l\le C[u]_{{s,q}}^{\xi_{l}}|u|_p^{1-\xi_{l}},\quad 
\forall u\in W^{s,q}(\mathbb {R}^N)\cap L^{p}(\mathbb {R}^N)
\end{equation}
for a positive constant $C>0$, where $\xi_{l}=\frac{Nq(l-p)}{l[Nq-p(N-sq)]}$,
$l\in(p,q_{s}^{*})$, $q_{s}^{*}=\frac{Nq}{N-sq}$. Hence,
\begin{equation}\label{eqn-(2.3)}
\begin{aligned}
J_\mu(u)
&\geq \frac{1}{p}\iint_{\mathbb{R}^{2N}}\frac{|u(x)-u(y)|^p}{|x-y|^{N+sp}}
 \mathrm{d}x\mathrm{d}y
 +\frac{1}{q}\iint_{\mathbb{R}^{2N}}\frac{|u(x)-u(y)|^q}{|x-y|^{N+sq}}\mathrm{d}x
 \mathrm{d}y\\
&\quad -CC_1a^{(1-\xi_{m_1})m_1}
 \Big(\iint_{\mathbb{R}^{2N}}\frac{|u(x)-u(y)|^q}{|x-y|^{N+sq}}\mathrm{d}x\mathrm{d}y
 \Big)^{\frac{m_1\xi_{m_1}}{q}}\\
&\quad -CC_2a^{(1-\xi_{m_2})m_{2}}
 \big(\iint_{\mathbb{R}^{2N}}\frac{|u(x)-u(y)|^q}{|x-y|^{N+sq}}\mathrm{d}x\mathrm{d}y
 \Big)^{\frac{m_2\xi_{m_2}}{q}}\\
&\quad -CC_3a^{(1-\xi_{q})q}
 \Big(\iint_{\mathbb{R}^{2N}}\frac{|u(x)-u(y)|^q}{|x-y|^{N+sq}}\mathrm{d}x\mathrm{d}y
 \Big)^{\frac{q\xi_{q}}{q}}\\
&\quad -CC_3a^{(1-\xi_{\beta})\beta}
 \Big(\iint_{\mathbb{R}^{2N}}\frac{|u(x)-u(y)|^q}{|x-y|^{N+sq}}\mathrm{d}x\mathrm{d}y
 \big)^{\frac{\beta\xi_{\beta}}{q}},
\end{aligned}
\end{equation}
as $m_1\in(p,\min\{p+\frac{sp^2}{N},m^{*}\})$ and 
$~m_2,~q,~\beta\in(p,\min\{q+\frac{spq}{N},m^{*}\})$. 
Clearly $m_1\xi_{m_1}$, $m_2\xi_{m_2}$, $q\xi_{q}$, $\beta\xi_{\beta}<q$. 
This ensures that $J_\mu$ is bounded from below.
\end{proof}

The above lemma ensures that $M_\mu(a):=\underset{u\in S(a)}\inf J_\mu(u)$ is 
well-defined.

\begin{lemma}\label{lem:2.2} 
Let $\mu,~a>0$, then $M_\mu(a)<0$.
\end{lemma}

\begin{proof}
By {\rm (A1)}, we infer that $\lim_{t\to 0}\frac{m_1F(t)}{t^{m_1}}=\alpha_1>0$,
which implies that, for some $\delta>0$,
\begin{equation}
\label{eqn-(2.4)}
\frac{m_1F(t)}{t^{m_1}}\ge\frac{\alpha_1}{2}
\end{equation}
for all $t\in(0,\delta]$.
From (A8), we can introduce the existence of a constant $C_3$ such that
\begin{equation}\label{eqn-(2.5)}
C_{3}\min\{|t|^{q},|t|^{\beta}\}\le G(t)\le C_{3}(|t|^{q}+|t|^{\beta}),\quad 
\forall t\in\mathbb{R}.
\end{equation}
Let $0<u_0\in S(a)\cap L^{\infty}(\mathbb{R}^N)$, we set
$$
H(u_0,r)(x)=e^{\frac{Nr}{p}}u_0(e^{r}x),\quad\forall x \in \mathbb{R}^N,\;
\forall r\in\mathbb{R}.
$$
It is easy to see that
$$
\int_{\mathbb{R}^{N}}|H(u_0,r)(x)|^{p}\mathrm{d}x=a^{p}.
$$
Furthermore, a straightforward calculation leads to
\begin{gather*}
\int_{\mathbb{R}^{N}} F(H(u_0,r)(x))\mathrm{d}x
=e^{-Nr}\int_{\mathbb{R}^{N}}F(e^{\frac{Nr}{p}}u_0(x))\mathrm{d}x, \\
\int_{\mathbb{R}^{N}} G(H(u_0,r)(x))\mathrm{d}x
=e^{-Nr}\int_{\mathbb{R}^{N}}G(e^{\frac{Nr}{p}}u_0(x))\mathrm{d}x.
\end{gather*}
Then, for $r<0$ and $|r|$ big enough, we have
$$0< e^{\frac{Nr}{p}}u_0(x)\le\delta,\quad\forall x\in\mathbb{R}^N.$$
Moreover, by \eqref{eqn-(2.4)}, we obtain
$$\int_{\mathbb{R}^{N}}F(H(u_0,r)(x))\mathrm{d}x\ge\frac{\alpha_1}{2m_1}e^{\frac{Nr(m_1-p)}{p}}\int_{\mathbb{R}^N}|u_0(x)|^{m_1}\mathrm{d}x,$$
and by \eqref{eqn-(2.5)}, we have
$$
\int_{\mathbb{R}^{N}}G(H(u_0,r)(x))\mathrm{d}x
\ge e^{-Nr}C_{3}\min\{e^{\frac{Nrq}{p}}
\int_{\mathbb{R}}|u_0(x)|^{q}\mathrm{d}x,e^{\frac{Nr\beta}{p}}
\int_{\mathbb{R}}|u_0(x)|^{\beta}\mathrm{d}x\}.
$$
Note that
\begin{align*}
\iint_{\mathbb{R}^{2N}}\frac{|H(u_0,r)(x)-H(u_0,r)(y)|^p}{|x-y|^{N+sp}}
 \mathrm{d}x\mathrm{d}y
&=e^{Nr}\iint_{\mathbb{R}^{2N}}\frac{|u_0(e^{r}x)-u_0(e^{r}y)|^{p}}{|x-y|^{N+sp}}
 \mathrm{d}x\mathrm{d}y\\
&=e^{rsp}\iint_{\mathbb{R}^{2N}}\frac{|u_0(x)-u_0(y)|^{p}}{|x-y|^{N+sp}}
 \mathrm{d}x\mathrm{d}y
\end{align*}
and
\begin{align*}
\iint_{\mathbb{R}^{2N}}\frac{|H(u_0,r)(x)-H(u_0,r)(y)|^q}{|x-y|^{N+sq}}
 \mathrm{d}x\mathrm{d}y
&=e^{Nr}\iint_{\mathbb{R}^{2N}}\frac{|u_0(e^{r}x)-u_0(e^{r}y)|^{q}}{|x-y|^{N+sq}}
 \mathrm{d}x\mathrm{d}y\\
&=e^{rsq}\iint_{\mathbb{R}^{2N}}\frac{|u_0(x)-u_0(y)|^{q}}{|x-y|^{N+sq}}
 \mathrm{d}x\mathrm{d}y.
\end{align*}
So,
\begin{align*}
J_\mu(H(u_0,r))
&\le\frac{e^{rps}}{p}\iint_{\mathbb{R}^{2N}}\frac{|u_0(x)-u_0(y)|^{p}}{|x-y|^{N+sp}}
 \mathrm{d}x\mathrm{d}y\\
&\quad +\frac{e^{rps}}{q}\iint_{\mathbb{R}^{2N}}\frac{|u_0(x)-u_0(y)|^{q}}{|x-y|^{N+sq}}
 \mathrm{d}x\mathrm{d}y\\
&\quad -\frac{\mu\alpha_1}{2m_1}e^{\frac{Nr(m_1-p)}{p}}\int_{\mathbb{R}^N}|u_0(x)|^{m_1}
 \mathrm{d}x\\
&\quad -e^{-Nr}C_{3}\min\{e^{\frac{Nrq}{p}}\int_{\mathbb{R}}|u_0(x)|^{q}
 \mathrm{d}x,e^{\frac{Nr\beta}{p}}\int_{\mathbb{R}}|u_0(x)|^{\beta}\mathrm{d}x\}.
\end{align*}
Since $0<s< 1$ , $m_1\in(p,\min\{p+\frac{sp^2}{N},m^{*}\})$ and 
$m_2,\alpha,\beta\in(p,\min\{q+\frac{spq}{N},m^{*}\})$, $r $ is less than 
$0$ and $|r|$ is large enough, we have
\begin{align*}
&\frac{e^{rps}}{p}\iint_{\mathbb{R}^{2N}}\frac{|u_0(x)-u_0(y)|^{p}}{|x-y|^{N+sp}}
 \mathrm{d}x\mathrm{d}y+\frac{e^{rps}}{q}\iint_{\mathbb{R}^{2N}}
 \frac{|u_0(x)-u_0(y)|^{q}}{|x-y|^{N+sq}}\mathrm{d}x\mathrm{d}y\\
&-\frac{\mu\alpha_1}{2m_1}e^{\frac{Nr(m_1-p)}{p}}\int_{\mathbb{R}^N}|u_0(x)|^{m_1}
 \mathrm{d}x-e^{-Nr}C_{3}\min\{e^{\frac{Nr\alpha}{p}}\int_{\mathbb{R}}|u_0(x)|^{\alpha}
 \mathrm{d}x,e^{\frac{Nr\beta}{p}}\int_{\mathbb{R}}|u_0(x)|^{\beta}\mathrm{d}x\}\\
&= A_r<0.
\end{align*}
Then
$J_\mu(H(u_0,r))\le A_r<0$.
Hence $M_\mu (a)<0$.
\end{proof}

\begin{lemma}\label{lem:2.3} 
If $\mu>0,a>0$, then
\begin{itemize}
\item[(i)] $a\mapsto M_\mu (a)$ is a continuous mapping;
\item[(ii)] if $a_1\in(0,a)$ and $a_2=(a^{p}-a_1^{p})^{1/p}$,
then $M_\mu(a)<M_\mu(a_1)+M_\mu(a_2)$.
\end{itemize}
\end{lemma}

\begin{proof}
(i) Let $a>0$ and $\{a_n\}\subset(0,+\infty)$ such that $a_n\to a$.
From the definition of  $M_\mu$, there exists $u_n\in S(a_n)$ such that 
$M_\mu(a_n)\le J_\mu(u_n)<M_\mu(a_n)+\frac{1}{n}$ for every $n\in\mathbb{N}^{+}$.
By Lemma \ref{lem:2.2},  $M_\mu(a_n)<0$.
Moreover, from \eqref{eqn-(2.3)}, $\{u_n\}$ is bounded in $X$.

Now consider $v_n:=\frac{a}{a_n}u_n\in S(a)$. From the boundedness of 
$\{u_n\}$ and $a_n\to a$, we have
\begin{align*}
M_\mu(a)&\leq  J_\mu(v_n)\\
&= \frac{1}{p}(\frac{a^{p}}{a_n^{p}}-1)\iint_{\mathbb{R}^{2N}}
\frac{|u_n(x)-u_n(y)|^{p}}{|x-y|^{N+sp}}\mathrm{d}x\mathrm{d}y
+\frac{1}{q}(\frac{a^{q}}{a_n^{q}}-1)\iint_{\mathbb{R}^{2N}}
\frac{|u_n(x)-u_n(y)|^{q}}{|x-y|^{N+sq}}\mathrm{d}x\mathrm{d}y\\
&\quad +\int_{\mathbb{R}^N}(\mu F(u_n)-\mu F(\frac{a}{a_n}u_n))\mathrm{d}x
+\int_{\mathbb{R}^N}( G(u_n)- G(\frac{a}{a_n}u_n))\mathrm{d}x+J_\mu(u_n)\\
&= J_\mu(u_n)+o_n(1).
\end{align*}
Let $n\to +\infty$, we can get $M_\mu(a)\le \underset{n\to +\infty}\liminf  M_\mu(a_n)$.
Let $\{w_n\}$ be a bounded minimizing sequence of $M_\mu(a)$ and 
$z_n:=\frac{a_n}{a}w_n\in S(a)$, repeat the above process we have
$$ 
M_\mu(a_n)\le J_\mu(z_n)=J_\mu(w_n)+o_n(1)\implies 
 \limsup_{n\to+\infty}  M_\mu(a_n)\le M_\mu(a),
$$
so we have $M_\mu(a_n)\to M_\mu(a)$.

(ii) For any fixed $a_1\in(0,a)$, we prove that
\begin{equation}\label{eqn-(2.6)}
M_\mu(\theta a_1)<{\theta}^{p}M_\mu(a_1),\quad\forall\theta>1.
\end{equation}
Indeed, if $\{u_n\}\in S(a_1)$ be a minimizing sequence for $M_\mu(a_1)$, 
then $u_n(\theta^{-\frac{p}{N}}x)\in S(\theta a_1)$.
Since $\theta>1$ and $\frac{p(N-sq)}{N}<\frac{p(N-sp)}{N}<p$, we have
\begin{align*}
M_\mu(\theta a_1)-\theta^{p}J_\mu(u_n)&\leq  J_\mu(u_n(\theta^{-\frac{p}{N}}x))
 -\theta^{p}J_\mu(u_n)\\
&= \frac{\theta^{\frac{p(N-sp)}{N}}-\theta^{p}}{p}
 \iint_{\mathbb{R}^{2N}}\frac{|u_n(x)-u_n(y)|^{p}}{|x-y|^{N+sp}}\mathrm{d}x\mathrm{d}y\\
&\quad +\frac{\theta^{\frac{p(N-sq)}{N}}-\theta^{p}}{q}
 \iint_{\mathbb{R}^{2N}}\frac{|u_n(x)-u_n(y)|^{q}}{|x-y|^{N+sq}}\mathrm{d}x\mathrm{d}y
\leq 0.
\end{align*}
As a consequence, $M_\mu(\theta a_1)\le{\theta}^{p}M_\mu(a_1)$.

If $M_\mu(\theta a_1)={\theta}^{p}M_\mu(a_1)$, then $[u_n]_{s,p}^{p}\to 0$ and 
$[u_n]_{s,q}^{q}\to 0$ as $n\to +\infty$,
which indicate that $\int_{\mathbb{R}^N}F(u_n)\mathrm{d}x\to 0 $ and 
$\int_{\mathbb{R}^N}G(u_n)\mathrm{d}x\to 0 $ by inequality \eqref{eqn-(2.2)}. Hence
\begin{align*}
0&>M_\mu(a_1)\\
&= \lim\limits_{n\to +\infty}J_\mu(u_n)\\
&= \frac{1}{p}\lim\limits_{n\to +\infty}\iint_{\mathbb{R}^{2N}}\frac{|u_n(x)-u_n(y)|^{p}}{|x-y|^{N+sp}}\mathrm{d}x\mathrm{d}y+\frac{1}{q}\lim\limits_{n\to +\infty}\iint_{\mathbb{R}^{2N}}\frac{|u_n(x)-u_n(y)|^{q}}{|x-y|^{N+sq}}\mathrm{d}x\mathrm{d}y\\
&\quad -\lim\limits_{n\to +\infty}\int_{\mathbb{R}^N}\mu F(u_n)\mathrm{d}x
 -\lim\limits_{n\to +\infty}\int_{\mathbb{R}^N} G(u_n)\mathrm{d}x
= 0,
\end{align*}
which is a contradiction. So we have $M_\mu(\theta a_1)<{\theta}^{p}M_\mu(a_1)$.
Perform the same steps again, we can get
\begin{equation}\label{eqn-(2.7)}
M_\mu(\theta a_2)<{\theta}^{p}M_\mu(a_2),\quad\forall\theta>1.
\end{equation}
Finally, applying \eqref{eqn-(2.6)} with $\theta=\frac{a}{a_1}>1$ 
and \eqref{eqn-(2.7)} with $\theta=\frac{a}{a_2}>1$ respectively, we obtain
$$
M_\mu(a)=\frac{a_1^{p}}{a^{p}}M_\mu(\frac{a}{a_1}a_1)
+\frac{a_2^{p}}{a^{p}}M_\mu(\frac{a}{a_2}a_2)<M_\mu(a_1)+M_\mu(a_2).
$$
\end{proof}

\begin{lemma}\label{pro:2.1}
Suppose that $\{u_n\}\subset S(a)$ is a minimizing sequence of $M_\mu(a)$. 
Then  for some subsequence either
\begin{itemize}
\item[(i)] $\{u_n\}$ is strongly convergent,
or
\item[(ii)] there exists a sequence $v_n(\cdot)=u_n(\cdot+y_n)$ with 
$|y_n|\to+\infty$ and $(y_n)\subset\mathbb{R}^N$, which is strongly 
convergent to a function $v\in S(a)$ with $J_\mu(v)=M_\mu(a)$.
\end{itemize}
\end{lemma}

\begin{proof}
By Lemma \ref{lem:2.1} it is easy to see that $\{u_n\}$ is bounded. 
Then there exists a subsequence $u_n\rightharpoonup u$ in $X$, which is still 
denoted as itself.
Assume $u\neq 0$ and $|u|_p=b$,
we can infer that $b\in(0,a]$. For the case of $b\in(0,a)$,
by fractional Br\'{e}zis-Lieb lemma \cite{ref1}, we have
$$
[u_n-u]_{s,p}^{p}=[u_n]_{s,p}^{p}-[u]_{s,p}^{p}+o_n(1)
$$
and
$$
[u_n-u]_{s,q}^{q}=[u_n]_{s,q}^{q}-[u]_{s,q}^{q}+o_n(1).
$$
Furthermore, under the assumption of $f$ and $g$, we obtain
$$
\int_{\mathbb{R}^N}F(u_n)\mathrm{d}x
=\int_{\mathbb{R}^N}F(u)\mathrm{d}x+\int_{\mathbb{R}^N}F(u_n-u)\mathrm{d}x+o_n(1)
$$
and
$$
\int_{\mathbb{R}^N}G(u_n)\mathrm{d}x=\int_{\mathbb{R}^N}G(u)\mathrm{d}x
+\int_{\mathbb{R}^N}G(u_n-u)\mathrm{d}x+o_n(1).
$$
Let $v_n=u_n-u$ and $|v_n|_p=d_n\to d$, we have $a^{p}=b^{p}+d^{p}$ and 
$d_n\in(0,a)$ for $n$ big enough. So,
$$
M_\mu(a)+o_n(1)=J_\mu(u_n)=J_\mu(u)+J_\mu(v_n)+o_n(1)\ge M_\mu(d_n)+M_\mu(b)+o_n(1).
$$
By the continuity of $a\mapsto M_\mu (a)$, we obtain
$$
M_\mu(a)\ge M_\mu(d)+M_\mu(b).
$$
This contradicts the conclusion of Lemma \ref{lem:2.3} (ii),
 where $a^{p}=b^{p}+d^{p}$.
Hence $|u|_p=a$.

Combining with $|u_n|_p=|u|_p=a$ and $u_n\rightharpoonup u$ in $L^{p}(\mathbb{R}^N)$, 
we obtain
\begin{equation}\label{eqn-(2.8)}
u_n\to u \quad \text{in} \quad L^{p}(\mathbb{R}^N).
\end{equation}
Combining inequality \eqref{eqn-(2.2)}, (A1) and (A2), we have
\begin{equation}\label{eqn-(2.9)}
\int_{\mathbb{R}^N}F(u_n)\mathrm{d}x\to\int_{\mathbb{R}^N}F(u)\mathrm{d}x.
\end{equation}
Similarly,
\begin{equation}\label{eqn-(2.10)}
\int_{\mathbb{R}^N}G(u_n)\mathrm{d}x\to\int_{\mathbb{R}^N}G(u)\mathrm{d}x.
\end{equation}
So
\begin{align*}
M_\mu(a)
&=J_\mu(u_n)+o_n(1)=J_\mu(u)+J_\mu(v_n)+o_n(1)\\
&\ge\frac{1}{p}[v_n]_{s,p}^{p}+\frac{1}{q}[v_n]_{s,q}^{q}+M_\mu(a)+o_n(1).
\end{align*}
This implies that $[v_n]_{s,p}^{p},[v_n]_{s,q}^{q}\le o_n(1)$.
So we have $v_n\to 0$ in $X$, which means $u_n\to u$ in $X$.

Assume $u=0$, i.e. $u_n\rightharpoonup 0$ in $X$.
Then, for some $\zeta_1,r_1>0$ and ${y_n}\subset\mathbb{R}^N$, we have
\begin{equation}\label{eqn-(2.11)}
\int_{B_{r_1}(y_n)}|u_n|^{p}\mathrm{d}x\ge\zeta_1, \quad \forall y_n\in\mathbb{R}^N.
\end{equation}
Otherwise by \cite[Lemma 2.1]{ref1} we have $u_n\to0$ in 
$L^{k}(\mathbb{R}^N)$ for all $k\in(p,p_{s}^{*})\cup (q,q_{s}^{*})$. 
Let $r_1\in (p, q_{s}^{*})$, and taking $p_1\in(p,p_{s}^{*})$ with $p<r_1$ and 
$q_1\in(q,q_{s}^{*})$ with $q_1>r_1$. Now let $t\in(0,1)$ be such that 
$r_1=tp_1+(1-t)q_1$, by the H\"{o}lder inequality we can write
$$
\int_{\mathbb{R}^N}|u_n|^{r_1}\mathrm{d}x
 =\int_{\mathbb{R}^N}|u_n|^{tp_1}|u_n|^{(1-t)q_1}\mathrm{d}x
 \le\Big(\int_{\mathbb{R}^N}|u_n|^{p_1}\mathrm{d}x\Big)^{t}
 \Big(\int_{\mathbb{R}^N}|u_n|^{q_1}\mathrm{d}x\Big)^{(1-t)}\to 0.
$$
Hence $u_n\to 0$ in $L^{r_1}(\mathbb{R}^N),~\forall r_1\in(p,q_{s}^{*})$,
which implies $F(u_n)\to 0$ and $G(u_n)\to 0$ in $L^{1}(\mathbb{R}^N)$.
However, this contradicts the fact that 
$$
0>M_\mu(a)+o_n(1)=J_\mu(u_n)\ge-\int_{\mathbb{R}^N}F(u_n)\mathrm{d}x
-\int_{\mathbb{R}^N}G(u_n)\mathrm{d}x.
$$
Hence \eqref{eqn-(2.11)} holds. Combining with $u=0$, the inequality \eqref{eqn-(2.11)} 
and Sobolev embedding, we deduce that $\{y_n\}$ is unbounded. 
Consider $v_n (\cdot)=u_n(x+y_n)$, it is easy to verify that $\{v_n\}$ is also 
a minimizing sequence of $M_\mu(a)$ and $\{v_n\}\subset S(a)$. As a result, 
it holds $v_n\rightharpoonup v$ in $X$, where $v\in X\backslash \{0\}$.
By an argument analogous to that used above, we can deduce $v_n\to v$ in $X$.
\end{proof}

\begin{lemma}\label{lem:2.5}
Assume {\rm (A1)--(A3), (A7)--(A9)} hold with $\mu>0$. 
Then, problem \eqref{eqn-(2.1)} has a positive radial solution $u$ and $\lambda<0$.
\end{lemma}

\begin{proof}
According to the definition of $J_\mu(u)$, we have $J_\mu(|u|)\le J_\mu(u)$.
In addition, we can also get $|u|\in S(a)$. Then, we deduce that
$$
M_\mu(a)=J_\mu(u)\ge J_\mu(|u|)\ge M_\mu(a),
$$
and we obtain $J_\mu(|u|)=M_\mu(a)$. Thus, instead of $|u|$, we can use $u$. 
If $u^{*}$ is the fractional Schwarz's Symmetrization of $u$ 
(\cite[Section 9.2]{ref5}), we obtain
\begin{gather*}
\iint_{\mathbb{R}^{2N}}\frac{|u(x)-u(y)|^{p}}{|x-y|^{N+sp}}
\mathrm{d}x\mathrm{d}y\ge \iint_{\mathbb{R}^{2N}}\frac{|u^{*}(x)
-u^{*}(y)|^{p}}{|x-y|^{N+sp}}\mathrm{d}x\mathrm{d}y,\\
\iint_{\mathbb{R}^{2N}}\frac{|u(x)-u(y)|^{q}}{|x-y|^{N+sq}}\mathrm{d}x\mathrm{d}y
\ge \iint_{\mathbb{R}^{2N}}\frac{|u^{*}(x)-u^{*}(y)|^{q}}{|x-y|^{N+sq}}
\mathrm{d}x\mathrm{d}y,
\end{gather*}
and by \cite[section 3.3]{ref22},
\begin{gather*}
\int_{\mathbb{R}^N} |u|^{p}\mathrm{d}x=\int_{\mathbb{R}^N} |u^{*}|^{p}\mathrm{d}x,\\
\int_{\mathbb{R}^N} G(u)\mathrm{d}x=\int_{\mathbb{R}^N} G(u^{*})\mathrm{d}x,\\
\int_{\mathbb{R}^N}\mu F(u)\mathrm{d}x=\int_{\mathbb{R}^N}\mu F(u^{*})\mathrm{d}x.
\end{gather*}
It is easy to prove that $u^{*}\in S(a)$ and $M_\mu(a)=J_\mu(u^{*})$. 
Thus, we replace $u$ by $u^{*}$.

From \cite{n4},  $u \in C^{\alpha}(\mathbb{R}^N)$ for some $\alpha\in (0,1)$. 
Now we show that $u(x)$ is positive for all $x\in\mathbb{R}^N$.
Assume by contradiction that there exists some $x_0\in\mathbb{R}^N$ satisfying 
$u(x_0)=0$, and there exists $x_1\in\mathbb{R}^N$ satisfying $u(x_1)>0$ since 
$u\neq 0$.
Thus, we can find a ball of sufficiently large radius
$R_1 > 0$ such that $x_0, x_1 \in B_{R_1/2}(0)$.
Then, combining this with the weak Harnack inequality \cite{n1}, we deduce that 
there exists a constant $C_8 > 0$ such that
$$
\sup_{y\in B_{R_1/2}(0)}  u(y)\le C_8\underset{y\in B_{R_1/2}(0)}\inf u(y),
$$
which contradicts the fact that
$$
\sup_{y\in B_{R_1/2}(0)}u(y)\ge u(x_1)>0\quad \text{and}\quad
\inf_{y\in B_{R_1/2}(0)}u(y)=u(x_0)=0.
$$
This implies that there is a positive radial solution $u$.

Next, we show that $\lambda<0$. By Lemma \ref{lem:2.2} we can assume that there 
exists a bounded minimizing sequence $\{u_n\}\in S(a)$.
Hence, we can obtain that there exists a constant $\lambda\in\mathbb{R}$ 
such that
\begin{equation}\label{eqn-(2.12)}
J_\mu'(u)=\lambda\Psi'(u)\quad  \text{in} \quad X',
\end{equation}
where $\Psi(u):=\int_{\mathbb{R}^N}|u|^{p}\mathrm{d}x$.
Then, from \eqref{eqn-(2.12)},
$$
(-\Delta)_p^{s} u+(-\Delta)_q^{s} u=\lambda |u|^{p-2}u+\mu f(u)+g(u),\quad
 x\in\mathbb{R}^N
$$
and
\begin{align*}
&\iint_{\mathbb{R}^{2N}}\frac{|u(x)-u(y)|^{p}}{|x-y|^{N+sp}}\mathrm{d}x\mathrm{d}y+\iint_{\mathbb{R}^{2N}}\frac{|u(x)-u(y)|^{q}}{|x-y|^{N+sq}}\mathrm{d}x\mathrm{d}y-\lambda \int_{\mathbb{R}^N} |u|^{p}\mathrm{d}x\\
&-\mu \int_{\mathbb{R}^N}f(u)u\mathrm{d}x-\int_{\mathbb{R}^N} g(u)u\mathrm{d}x=0.
\end{align*}
By (A3), (A8), and  that $M_\mu(a)=J_\mu(u)<0$,  we have
\begin{align*}
0&> J_\mu(u)-\frac{1}{q}(\iint_{\mathbb{R}^{2N}}
 \frac{|u(x)-u(y)|^{p}}{|x-y|^{N+sp}}\mathrm{d}x\mathrm{d}y
 +\iint_{\mathbb{R}^{2N}}\frac{|u(x)-u(y)|^{q}}{|x-y|^{N+sq}}\mathrm{d}x
 \mathrm{d}y-\lambda\int_{\mathbb{R}^N} |u|^{p}\mathrm{d}x\\
&\quad -\mu\int_{\mathbb{R}^N} f(u)u\mathrm{d}x-\int_{\mathbb{R}^N}g(u)u\mathrm{d}x)\\
&= (\frac{1}{p}-\frac{1}{q})\iint_{\mathbb{R}^{2N}}\frac{|u(x)-u(y)|^{p}}{|x-y|^{N+sp}}
 \mathrm{d}x\mathrm{d}y+\frac{\lambda}{q}\int_{\mathbb{R}^N} |u|^{p}\mathrm{d}x
 +\frac{\mu}{q} \int_{\mathbb{R}^N} f(u)u\mathrm{d}x-\mu\int_{\mathbb{R}^N}
 F(u)\mathrm{d}x\\
&\quad +\frac{1}{q}\int_{\mathbb{R}^N}g(u)u\mathrm{d}x-\int_{\mathbb{R}^N}
 G(u)\mathrm{d}x\\
&\geq \frac{1}{q}\int_{\mathbb{R}^N}\lambda |u|^{p}\mathrm{d}x,
\end{align*}
this implies that $\lambda<0$.
\end{proof}

From Lemma \ref{lem:2.5},  we obtain the following lemma.

\begin{lemma}\label{cor:2.1} 
Fix $a>0$ and let $0<\mu_1<\mu_2$.
Then, $M_{\mu_2(a)}<M_{\mu_1}(a)<0$.
\end{lemma}

\section{Non-autonomous case}

We will state some properties of the functional $I_{\varepsilon}$,
which is restricted to $S(a)$.
Firstly, we define $I_{\rm max},I_\infty:X\to \mathbb{R}$ as
\begin{align*}
I_{\rm max}(u)
&= \frac{1}{p}\iint_{\mathbb{R}^{2N}}\frac{|u(x)
 -u(y)|^{p}}{|x-y|^{N+sp}}\mathrm{d}x\mathrm{d}y
  +\frac{1}{q}\iint_{\mathbb{R}^{2N}}\frac{|u(x)-u(y)|^{q}}{|x-y|^{N+sq}}\mathrm{d}x\mathrm{d}y\\
&\quad -\int_{\mathbb{R}^N}h_{\rm max}F(u)\mathrm{d}x
 -\int_{\mathbb{R}^N}G(u)\mathrm{d}x
\end{align*}
and
\begin{align*}
I_\infty(u)
&= \frac{1}{p}\iint_{\mathbb{R}^{2N}}\frac{|u(x)-u(y)|^{p}}
 {|x-y|^{N+sp}}\mathrm{d}x\mathrm{d}y+\frac{1}{q}
 \iint_{\mathbb{R}^{2N}}\frac{|u(x)-u(y)|^{q}}{|x-y|^{N+sq}}\mathrm{d}x\mathrm{d}y\\
&\quad -\int_{\mathbb{R}^N}h_{\infty}F(u)\mathrm{d}x
-\int_{\mathbb{R}^N}G(u)\mathrm{d}x.
\end{align*}
Moreover, by Lemma \ref{lem:2.2} we know that
$$ 
M_\infty(a)=\inf_{u\in S(a)}I_\infty(u),\quad
M_\varepsilon(a)=\inf_{u\in S(a)}I_\varepsilon(u),\quad
M_{\rm max}(a)=\inf_{u\in S(a)}I_{\rm max}(u).
$$
Then, by Lemma \ref{cor:2.1} and $h_\infty<h_{\rm max}$, we can immediately obtain
\begin{equation}\label{eqn-(3.1)}
M_{\rm max}(a)<M_\infty(a)<0.
\end{equation}
Now, we fix $0<\rho_1=\frac{1}{p}(M_\infty(a)-M_{\rm max}(a))$.

\begin{lemma}\label{lem:3.1}
 $\lim_{\varepsilon\to 0^{+}}  M_\varepsilon(a)\le M_{\rm max}(a)$, and there exists 
 $\varepsilon_0>0$ such that $M_\varepsilon(a)<M_\infty(a)$ for all 
 $0<\varepsilon<\varepsilon_0$.
\end{lemma}

\begin{proof} 
Let $u_0\in S(a)$ satisfying $I_{\rm max}(u_0)=M_{\rm max}(a)$.
A simple calculation leads to
\begin{align*}
M_\varepsilon(a)&\leq  I_\varepsilon(u_0)\\
&= \frac{1}{p}\iint_{\mathbb{R}^{2N}}\frac{|u_0(x)-u_0(y)|^{p}}{|x-y|^{N+sp}}
 \mathrm{d}x\mathrm{d}y+\frac{1}{q}\iint_{\mathbb{R}^{2N}}
 \frac{|u_0(x)-u_0(y)|^{q}}{|x-y|^{N+sq}}\mathrm{d}x\mathrm{d}y\\
&\quad -\int_{\mathbb{R}^N}h(\varepsilon x)F(u_0)\mathrm{d}x
 -\int_{\mathbb{R}^N}G(u_0)\mathrm{d}x.
\end{align*}
Letting $\varepsilon\to 0^{+}$ and applying (A6) we can obtain
$$
\limsup\limits_{\varepsilon\to 0^{+}} M_ \varepsilon (a)
\le\lim_{\varepsilon\to 0^{+}}I_\varepsilon(u_0)=I_{\rm max}(u_0)=M_{\rm max}(a).
$$
According to \eqref{eqn-(3.1)}, we have $M_{\varepsilon}(a)<M_{\infty}(a)$ for 
$\varepsilon$ small enough.
\end{proof}

The following two lemmas will be used to prove the $(PS)_c$ condition for 
$I_\varepsilon$ at certain levels.

\begin{lemma}\label{lem:3.2}
Assume that $\{u_n\}\subset S(a)$ is a minimizing sequence with 
$I_\varepsilon(u_n)\to c$ and $c < M_{\rm max}(a)+\rho_1 < M_{\infty}(a)$.
If $u_n\rightharpoonup u$ in $X$, then $u \neq  0$.
\end{lemma}

\begin{proof}
Firstly, we assume that the conclusion is not true, i.e. $u\equiv0$. 
Then, we obtain
$$
c=M_{\varepsilon}(a)=I_\varepsilon(u_n)+o_n(1)
=I_\infty(u_n)+\int_{\mathbb{R}^N}(h_\infty-h(\varepsilon x))F(u_n)\mathrm{d}x+o_n(1).
$$
From (A5), for any arbitrary number $\zeta_2>0$, there exists some constant $ R_2 > 0$ 
such that
$$
h_{\infty}\ge h(x)-\zeta_2,\quad|x|>R_2.
$$
Then
\begin{align*}
c&= I_\varepsilon(u_n)+o_n(1)\\
&\geq  I_\infty(u_n)+\int_{B_{R_2/\varepsilon}(0)}(h_\infty-
h(\varepsilon x))F(u_n)\mathrm{d}x-\zeta_2\int_{B_{R_2/\varepsilon}^{c}(0)}F(u_n)
\mathrm{d}x+o_n(1).
\end{align*}
Since $\{u_n\}$ is bounded in $X$, then for some constant $C_9 > 0$, we have
\begin{align*}
\int_{\mathbb{R}^N}F(u_n)\mathrm{d}x
&\le C_1a^{(1-\xi_{m_1})m_1}(\iint_{\mathbb{R}^{2N}}
 \frac{|u(x)-u(y)|^{q}}{|x-y|^{N+sq}}\mathrm{d}x\mathrm{d}y)^{\frac{m_1\xi_{m_1}}{q}}\\
&\quad +C_2a^{(1-\xi_{m_2})m_2}(\iint_{\mathbb{R}^{2N}}
 \frac{|u(x)-u(y)|^{q}}{|x-y|^{N+sq}}\mathrm{d}x\mathrm{d}y)^{\frac{m_2\xi_{m_2}}{q}}\\
&\le C_9.
\end{align*}
From the fact that $u_n\to 0$ in $L^{k_1}(B_{R/\varepsilon}(0))$ when 
$k_1\in[1,q_{s}^{*})$, we have
$$
c=I_\varepsilon(u_n)+o_n(1)\ge I_\infty(u_n)-\zeta_2 C_9+o_n(1)\ge M_\infty (a)
-\zeta_2 C_9+o_n(1).
$$
This together with small enough $\zeta_2> 0$, we obtain
$c\ge M_\infty (a)$.
This contradicts $c < M_{\rm max}(a) + \rho_1 < M_\infty(a)$. 
Therefore, we can obtain $u\neq 0$.
\end{proof}

\begin{lemma}\label{lem:3.3} 
Assume that $\{u_n\}\subset S(a)$  is a
$(PS)_c$ sequence of $I_\varepsilon$ satisfying $u_n\rightharpoonup u_\varepsilon$ 
in $X$ where $c < M_{\rm max}(a) + \rho_1 <M_{\infty}(a)$, that is, as $n\to+\infty$,
$$
I_\varepsilon(u_n)\to c \quad \text{and}\quad\|I_\varepsilon|_{S(a)}'(u_n)\|_{X'}\to 0.
$$
Then
$$
\liminf_{n\to+\infty}|u_n-u_\varepsilon|_p^{p}\ge \beta_1,
$$
where $u_n\not\to u_\varepsilon $ in $X$ and $\beta_1>0$ independent of 
$\varepsilon\in(0,\varepsilon_0)$.
\end{lemma}

\begin{proof}
Firstly, define the functional 
$\Gamma(u)=\frac{1}{p}\int_{\mathbb{R}^N}|u|^{p}\mathrm{d}x$. 
We can see that $S(a)=\Gamma^{-1}({a^{p}/p})$. 
According to \cite[Proposition 5.12]{ref31}, there exists 
$\{\lambda_n\}\subset\mathbb{R}$ such that
$$
\|I_\varepsilon'(u_n)-\lambda_n\Gamma'(u_n)\|_{X'}\to 0
$$ 
as $n\to+\infty$.
Since $I_\varepsilon$ is bounded from below and $I_{\rm max}$ is forced, it follows that
$\{u_n\}$ is bounded in $X$. 
This ensures that
$\{\lambda_n\}$ is bounded, up to a subsequence, we may assume
$\lambda_n\to \lambda_\varepsilon$ as $n\to+\infty$.
Now we prove that
\begin{gather}\label{eqn-(3.2)}
I_\varepsilon'(u_\varepsilon)-\lambda_\varepsilon\Gamma'(u_\varepsilon)=0\quad 
\text{in } X',\\
\label{eqn-(3.3)}
\|I_\varepsilon'(v_n)-\lambda_\varepsilon\Gamma'(v_n)\|_{X'}\to 0\quad 
\text{as } n\to+\infty,
\end{gather}
where $v_n=u_n-u_\varepsilon$.
Since $\|I_\varepsilon'(u_n)-\lambda_n\Gamma'(u_n)\|_{X'}\to 0$,
to prove \eqref{eqn-(3.2)}, it is sufficient to prove that for all $\phi\in X$,
\begin{gather}\label{eqn-(3.5)}
\begin{aligned}
&\iint_{\mathbb{R}^{2N}}\frac{|u_n(x)-u_n(y)|^{p-2}(u_n(x)-u_n(y))(\phi(x)
 -\phi(y))}{|x-y|^{N+sp}}\mathrm{d}x\mathrm{d}y\\
&\to\iint_{\mathbb{R}^{2N}}\frac{|u_\varepsilon(x)
 -u_\varepsilon(y)|^{p-2}(u_\varepsilon(x)-u_\varepsilon(y))(\phi(x)
 -\phi(y))}{|x-y|^{N+sp}}\mathrm{d}x\mathrm{d}y,
\end{aligned} \\
\label{eqn-(3.6)}
\begin{aligned}
&\iint_{\mathbb{R}^{2N}}\frac{|u_n(x)-u_n(y)|^{q-2}(u_n(x)-u_n(y))(\phi(x)-\phi(y))}{|x-y|^{N+sq}}\mathrm{d}x\mathrm{d}y\\
&\to\iint_{\mathbb{R}^{2N}}\frac{|u_\varepsilon(x)-u_\varepsilon(y)|^{q-2}(u_\varepsilon(x)-u_\varepsilon(y))(\phi(x)-\phi(y))}{|x-y|^{N+sq}}\mathrm{d}x\mathrm{d}y,
\end{aligned} \\
\label{eqn-(3.7)}
\begin{aligned}
&\int_{\mathbb{R}^N}h(\varepsilon x)f(u_n)\phi\mathrm{d}x
\to\int_{\mathbb{R}^N}h(\varepsilon x)f(u_\varepsilon)\phi\mathrm{d}x,
\end{aligned} \\
\label{eqn-(3.8)}
\int_{\mathbb{R}^N}g(u_n)\phi\mathrm{d}x\to
 \int_{\mathbb{R}^N}g(u_\varepsilon)\phi\mathrm{d}x,\\
\label{eqn-(3.9)}
\lambda_n \int_{\mathbb{R}^N}|u_n|^{p-2}u_n\phi\mathrm{d}x
\to\lambda_\varepsilon \int_{\mathbb{R}^N}|u_\varepsilon|^{p-2}
u_\varepsilon\phi\mathrm{d}x.
\end{gather}

We first prove \eqref{eqn-(3.5)} holds using the technique in \cite{n5}.
From $u_n \rightharpoonup u_{\varepsilon}$ in $X$,  $u_n \to u_{\varepsilon}$ in 
$L_{\rm loc}^t(\mathbb{R}^N)$ for all $t \in [1,q_{s}^{\ast})$.
For any $\phi \in C_c^{\infty}(\mathbb{R}^N)$, define
\[
\Phi(x,y) = \frac{\phi(x) - \phi(y)}{|x - y|^{\frac{N+sp}{p}}} \in L^p(\mathbb{R}^{2N}).
\]
We also define
\begin{gather*}
U_n(x,y) = \frac{|u_n(x) - u_n(y)|^{p-2}(u_n(x) - u_n(y))}{|x - y|^{\frac{N+sp}{p'}}}
\in L^{p'}(\mathbb{R}^{2N}), \\
U_\varepsilon(x,y) = \frac{|u_\varepsilon(x) 
- u_\varepsilon(y)|^{p-2}(u_\varepsilon(x) - u_\varepsilon(y))}{|x - y|^{\frac{N+sp}{p'}}}
\in L^{p'}(\mathbb{R}^{2N}),
\end{gather*}
where $p' = \frac{p}{p-1}$. It is easy to see that $\{U_n\}$ is  bounded in 
$L^{p'}(\mathbb{R}^{2N})$. Since $L^{p'}(\mathbb{R}^{2N})$ is reflexive,  
there exists a subsequence, still denoted by $\{U_n\}$, such
that  $U_n \rightharpoonup U_\varepsilon$ in $L^{p'}(\mathbb{R}^{2N})$, hence
\[
\iint_{\mathbb{R}^{2N}} U_n(x,y)\Phi(x,y)\mathrm{d}x\mathrm{d}y 
\to \iint_{\mathbb{R}^{2N}} U_\varepsilon(x,y)\Phi(x,y)\mathrm{d}x\mathrm{d}y.
\]
Consequently, we obtain \eqref{eqn-(3.5)}.
Using the same technique it can be shown that \eqref{eqn-(3.6)} holds.

Next we prove that \eqref{eqn-(3.7)} holds.
By the condition of $h$ we know that $h$ is bounded, then we only need to prove
\begin{equation}\label{eqn-(3.11)}
\int_{\mathbb{R}^N}(f(u_n)-f(u_\varepsilon))\phi\mathrm{d}x\to0.
\end{equation}
From conditions (A1) and (A2) we know that
$$
|f(t)|\leq C_1|t|^{m_1-1}+C_2|t|^{m_2-1},
$$
then there exists a constant $C_{10}$ such that
$$
|f(u_n)-f(u_\varepsilon)|\leq C_{10}(|u_n|^{m_1-1}+|u_n|^{m_2-1}
+|u_\varepsilon|^{m_1-1}+|u_\varepsilon|^{m_2-1}).
$$
To prove \eqref{eqn-(3.11)}, we divide the interval of integration into two parts, 
i.e.\ for a sufficiently large $R_3>0$ independent of $N$,
$$
\int_{\mathbb{R}^N}(f(u_n)-f(u_\varepsilon))\phi\mathrm{d}x
=\int_{B_{R_3}}(f(u_n)-f(u_\varepsilon))\phi\mathrm{d}x
 +\int_{\mathbb{R}^N\backslash{B_{R_3}}}(f(u_n)-f(u_\varepsilon))\phi\mathrm{d}x.
$$
(i) For integrals within $B_{R_3}$. Using H\"{o}lder's inequality we have
\begin{align*}
|\int_{B_{R_3}}(f(u_n)-f(u_\varepsilon))\phi\mathrm{d}x|
&\le\int_{B_{R_3}}|f(u_n)-f(u_\varepsilon)||\phi|\mathrm{d}x\\
&\le \big(\int_{B_{R_3}}|f(u_n)-f(u)|^{(q_{s}^{*})'}\Big)^{\frac{1}{(q_{s}*)'}}
 \Big(\int_{B_{R_3}}|\phi|^{q_{s}^{*}}\Big)^{\frac{1}{q_{s}^{*}}},
\end{align*}
where $\frac{1}{q_{s}^{*}}+\frac{1}{(q_{s}^{*})'}=1$. 
Because of $u_n\rightharpoonup u$ in $X_{R_3}=W^{s,p}(B_{R_3})\cap W^{s,q}(B_{R_3})$,
 and $X_{R_3}\hookrightarrow\hookrightarrow L^{k_2}(B_{R_3})$, 
 where $k_2\in[1,q_{s}^{*})$, it follows that $u_n\to u$ in $L^{k_2}(B_{R_3})$.
Also by the condition (A2) and the continuity of $f$ it follows that 
$f\in C(L^{k_2}(B_{R_3}),L^{(q_{s}^{*})'}(B_{R_3}))$, and 
$\int_{B_{R_3}}|f(u_n)-f(u_\varepsilon)|^{(q_{s}^{*})'}\mathrm{d}x\to 0$,
hence
$$
\int_{B_{R_3}}(f(u_n)-f(u_\varepsilon))\phi\mathrm{d}x\to 0.
$$

(ii) For integrals outside the $B_{R_3}$. By H\"{o}lder's  inequality we have
\begin{align*}
|\int_{\mathbb{R}^N\backslash{B_{R_3}}}(f(u_n)-f(u_\varepsilon))
\phi\mathrm{d}x|
&\leq \Big(\int_{\mathbb{R}^N\backslash{B_{R_3}}}
|f(u_n)-f(u_\varepsilon)|^{p'}\mathrm{d}x\Big)^{1/p'}
\Big(\int_{\mathbb{R}^N\backslash{B_{R_3}}}|\phi|^{p}\mathrm{d}x\Big)^{1/p}\\
&\leq C_{10}\Big[\Big(\int_{\mathbb{R}^N\backslash{B_{R_3}}}|u_n|^{(m_1-1)p'}
 \mathrm{d}x\Big)^{1/p'}\Big(\int_{\mathbb{R}^N\backslash{B_{R_3}}}|\phi|^{p}
 \mathrm{d}x\Big)^{1/p}\\
&\quad +\Big(\int_{\mathbb{R}^N\backslash{B_{R_3}}}|u_\varepsilon|^{(m_1-1)p'}
 \mathrm{d}x\Big)^{1/p'}\Big(\int_{\mathbb{R}^N\backslash{B_{R_3}}}|\phi|^{p}
 \mathrm{d}x\Big)^{1/p}\\
&\quad +\Big(\int_{\mathbb{R}^N\backslash{B_{R_3}}}|u_n|^{(m_2-1)p'}
 \mathrm{d}x\Big)^{1/p'}
 \Big(\int_{\mathbb{R}^N\backslash{B_{R_3}}}|\phi|^{p}\mathrm{d}x\Big)^{1/p}\\
&\quad +\Big(\int_{\mathbb{R}^N\backslash{B_{R_3}}}|u_\varepsilon|^{(m_2-1)p'}
 \mathrm{d}x\Big)^{1/p'}\Big(\int_{\mathbb{R}^N\backslash{B_{R_3}}}|\phi|^{p}
 \mathrm{d}x\Big)^{1/p}\Big],
\end{align*}
where $\frac{1}{p}+\frac{1}{p'}=1$. By (A1) we have $p<(m_1-1)p'<q_{s}^{*}$, then  
$\{u_n\}$ is bounded  in $L^{(m_1-1)p'}(\mathbb{R}^N)$.  Therefore,
there exists a large enough $R_3>0$, such that 
$$
\int_{\mathbb{R}^N\backslash{B_{R_3}}}|u_n|^{(m_1-1)p'}\mathrm{d}x<\varepsilon.
$$
From the definition of $\phi$ it follows that $\phi$ is bounded in $X$, i.e.\
there exists a constant $M>0$ such that $|\phi|_p<M$.
Hence,
$$
\Big(\int_{\mathbb{R}^N\backslash{B_{R_3}}}|u_n|^{(m_1-1)p'}
\mathrm{d}x\Big)^{1/p'} \Big(\int_{\mathbb{R}^N\backslash{B_{R_3}}}
|\phi|^{p}\mathrm{d}x\Big)^{1/p}<\varepsilon M. 
$$
For the same reason,
$$
\Big(\int_{\mathbb{R}^N\backslash{B_{R_3}}}|u_\varepsilon|^{(m_1-1)p'}\mathrm{d}x
 \Big)^{1/p'}
 \Big(\int_{\mathbb{R}^N\backslash{B_{R_3}}}|\phi|^{p}\mathrm{d}x\Big)^{1/p}
 <\varepsilon M.
 $$
By (A2) and repeat the above steps, we have
\begin{gather*}
\Big(\int_{\mathbb{R}^N\backslash{B_{R_3}}}|u_n|^{(m_2-1)p'}\mathrm{d}x\Big)^{1/p'}
\Big(\int_{\mathbb{R}^N\backslash{B_{R_3}}}|\phi|^{p}\mathrm{d}x\Big)^{1/p}
<\varepsilon M,\\
\Big(\int_{\mathbb{R}^N\backslash{B_{R_3}}}|u_\varepsilon|^{(m_2-1)p'}\mathrm{d}x
\Big)^{1/p'}
\Big(\int_{\mathbb{R}^N\backslash{B_{R_3}}}|\phi|^{p}\mathrm{d}x\Big)^{1/p}
<\varepsilon M.
\end{gather*}
Hence, $|\int_{\mathbb{R}^N\backslash{B_{R_3}}}(f(u_n)-f(u_\varepsilon))
\phi\mathrm{d}x|<4\varepsilon M$.
By the arbitrariness of $\varepsilon$, we have
$$
\int_{\mathbb{R}^N}(f(u_n)-f(u_\varepsilon))\phi\mathrm{d}x
=\int_{B_{R_3}}(f(u_n)-f(u_\varepsilon))\phi\mathrm{d}x
 +\int_{\mathbb{R}^N\backslash{B_{R_3}}}(f(u_n)-f(u_\varepsilon))\phi\mathrm{d}x\to 0.
$$
Summarizing, we have \eqref{eqn-(3.7)}.
Carry out the same method once more, we can obtain \eqref{eqn-(3.8)}.

To prove \eqref{eqn-(3.9)}, let $A(u_n)=|u_n|^{p-2}u_n$. It follows from the 
boundness of $\{u_n\}$ in $X$, there exists a constant $C_{11}$ such that
$$
|A(u_n)|_{L^{p'}}=||u_n|^{p-2}u_n|_{L^{p'}}=|u_n|_{L^p}^{p-1}\le C_{11}^{p-1}.
$$
Then for each $ \varepsilon_1>0$ there exists a compact support function $\psi_1$ 
such that for every $\phi\in X$ with $|\phi-\psi_1|_{L^p}<\varepsilon_1$,
we  obtain
\begin{align*}
\int_{\mathbb{R}^N}(A(u_n)-A(u_\varepsilon))\phi\mathrm{d}x
&\leq |A(u_n)-A(u_\varepsilon)|_{L^{p'}}|\phi-\psi_1|_{L^p}
 +\int_{\mathbb{R}^N}(A(u_n)-A(u_\varepsilon))\psi_1\mathrm{d}x\\
&<\varepsilon_1 C_{11}^{p-1}+\int_{\mathbb{R}^N}(A(u_n)-A(u_\varepsilon))
 \psi_1\mathrm{d}x.
\end{align*}
By the arbitrariness of $\varepsilon_1$, we have
$$
\int_{\mathbb{R}^N}(A(u_n)-A(u_\varepsilon))\phi\mathrm{d}x\to 0.
$$
Hence \eqref{eqn-(3.9)} holds.

From (A3), we have
\begin{align*}
0> M_{\rm max}(a) + \rho_1>c
&=\liminf_{n\to+\infty}I_\varepsilon(u_n)\\
&=\liminf_{n\to+\infty}(I_\varepsilon(u_n)
 -\frac{1}{q}\langle I_{\varepsilon}' (u_n),u_n\rangle+\frac{1}{q}\lambda_n a^{p})\\
&\ge\frac{1}{q}\lambda_\varepsilon a^{p},
\end{align*}
this implies
$$
\limsup_{\varepsilon\to0}\lambda_\varepsilon\le\frac{q(M_{\rm max}(a) 
+ \rho_1)}{a^{p}}<0.
$$
Then, there exists a constant $\lambda^{*}$ independent of $\varepsilon$
that satisfies $\lambda_\varepsilon<\lambda^{*}<0$, so
\begin{align*}
&\iint_{\mathbb{R}^{2N}}\frac{|v_n(x)-v_n(y)|^{p}}{|x-y|^{N+sp}}
 \mathrm{d}x\mathrm{d}y
 +\iint_{\mathbb{R}^{2N}}\frac{|v_n(x)-v_n(y)|^{q}}{|x-y|^{N+sq}}
 \mathrm{d}x\mathrm{d}y-\lambda_\varepsilon\int_{\mathbb{R}^N}| v_n|^{p}\mathrm{d}x\\
&=\int_{\mathbb{R}^N}h(\varepsilon x)f(v_n)v_n\mathrm{d}x
 +\int_{\mathbb{R}^N}g(v_n)v_n\mathrm{d}x+o_n(1)
\end{align*}
and
\begin{align*}
&\iint_{\mathbb{R}^{2N}}\frac{|v_n(x)-v_n(y)|^{p}}{|x-y|^{N+sp}}
 \mathrm{d}x\mathrm{d}y+\iint_{\mathbb{R}^{2N}}
 \frac{|v_n(x)-v_n(y)|^{q}}{|x-y|^{N+sq}}\mathrm{d}x\mathrm{d}y
 -\lambda^{*}\int_{\mathbb{R}^N}| v_n|^{p}\mathrm{d}x\\
&\le\int_{\mathbb{R}^N}h(\varepsilon x)f(v_n)v_n\mathrm{d}x
 +\int_{\mathbb{R}^N}g(v_n)v_n\mathrm{d}x+o_n(1).
\end{align*}

According to (A1) and (A2), we have 
$|f(t)|<C_1|t|^{m_1-1}+C_1|t|^{m_2-1}$ for all $t\in\mathbb{R}$. Then
\begin{align*}
\int_{\mathbb{R}^N}f(v_n)v_n\mathrm{d}x
&\le \int_{\mathbb{R}^N}|f(v_n)||v_n|\mathrm{d}x\\
&\le\int_{\mathbb{R}^N}(C_1|v_n|^{m_1-1}+C_2|v_n|^{m_2-1})|v_n|\mathrm{d}x\\
&\le C_1\int_{\mathbb{R}^N}|v_n|^{m_1}\mathrm{d}x
+C_2\int_{\mathbb{R}^N}|v_n|^{m_2}\mathrm{d}x\,.
\end{align*}
From (A8), we have $|g(t)|\le C_3(|t|^{q-1}+|t|^{\beta-1})$ for all $t\in\mathbb{R} $, 
hence
$$
\int_{\mathbb{R}^N}|g(v_n)v_n|\mathrm{d}x
 \le C_3\int_{\mathbb{R}^N}|v_n|^{q}+|v_n|^{\beta}\mathrm{d}x
 =C_3\int_{\mathbb{R}^N}|v_n|^{q}\mathrm{d}x
 +C_3\int_{\mathbb{R}^N}|v_n|^{\beta}\mathrm{d}x.
$$
So, we have
\begin{align*}
&\iint_{\mathbb{R}^{2N}}\frac{|v_n(x)-v_n(y)|^{p}}{|x-y|^{N+sp}}
 \mathrm{d}x\mathrm{d}y+\iint_{\mathbb{R}^{2N}}\frac{|v_n(x)
 -v_n(y)|^{q}}{|x-y|^{N+sq}}\mathrm{d}x\mathrm{d}y
 +C_0\int_{\mathbb{R}^N}| v_n|^{p}\mathrm{d}x\\
&\le h_{\rm max}\int_{\mathbb{R}^N}f(v_n)v_n\mathrm{d}x
 +\int_{\mathbb{R}^N}g(v_n)v_n\mathrm{d}x\\
&\le C_1h_{\rm max}\int_{\mathbb{R}^N}|v_n|^{m_1}\mathrm{d}x
 +C_2h_{\rm max}\int_{\mathbb{R}^N}|v_n|^{m_2}\mathrm{d}x
 +C_3\int_{\mathbb{R}^N}|v_n|^{q}\mathrm{d}x
 +C_3\int_{\mathbb{R}^N}|v_n|^{\beta}\mathrm{d}x+o_n(1)
\end{align*}
for some constant $C_0=-\lambda^{*}>0$ independent of 
$\varepsilon\in(0,\varepsilon_0)$. Since $v_n\not\to0$  in $X$, we can assume that
$\underset{n\to+\infty}\liminf \|v_n\|_X > C^{*} > 0$.
Thus, for some constant $C_{12}>0$, by \eqref{eqn-(2.2)}, we infer that
\begin{equation}\label{eqn-(3.12)}
\begin{aligned}
C_{12}
&\leq \underset{n\to+\infty}\liminf(|v_n|_{m_1}^{m_1}+|v_n|_{m_2}^{m_2}+|v_n|_{q}^{q}
 +|v_n|_{\beta}^{\beta})\\
&\le \underset{n\to+\infty}\liminf C_{13}([v_n]_{{s,q}}^{m_1\xi_{m_1}}
 |v_n|_{p}^{m_1(1-\xi_{m_1})}+[v_n]_{{s,q}}^{m_2\xi_{m_2}}|v_n|_{p}^{m_2(1-\xi_{m_2})}\\
&\quad +[v_n]_{{s,q}}^{q\xi_{q}}|v_n|_{p}^{q(1-\xi_{q})}
 +[v_n]_{{s,q}}^{\beta\xi_{\beta}}|v_n|_{p}^{\beta(1-\xi_{\beta})})\\
&\leq  \liminf_{n\to+\infty}  C_{13}K_1(|v_n|_{p}^{m_1(1-\xi_{m_1})}
 +|v_n|_{p}^{m_2(1-\xi_{m_2})}+|v_n|_{p}^{q(1-\xi_{q})}+|v_n|_{p}^{\beta(1-\xi_{\beta})})
\end{aligned}
\end{equation}
where $K_1>0$ is independent of $\varepsilon\in(0,\varepsilon_0)$ with 
$[v_n]_{{s,q}}\le K_1$ for all $n\in\mathbb{N}$.
Since $m_1\in(p,\min\{p+\frac{sp^2}{N},m^{*}\})$, 
$m_2, q,\beta\in(p,\min\{q+\frac{spq}{N},m^{*}\})$, then 
$m_1(1-\xi_{m_1}),~m_2(1-\xi_{m_2}),~q(1-\xi_{q}),~\beta(1-\xi_{\beta})<p$,
hence
\begin{equation}\label{eqn-(3.13)}
\liminf_{n\to +\infty} |v_n|_p^{p}\ge C_{15}
\end{equation}
for a constant $C_{15}>0$. This complete the proof.
\end{proof}

Next, we consider 
$0<\rho<\min\{{\frac{1}{p},\frac{\beta_1}{a^{p}}}\}(M_\infty(a)-M_{\rm max}(a))$.

\begin{lemma}\label{lem:3.4} 
Assume that $0<\varepsilon<\varepsilon_0$ and $c<M_{\rm max}(a)+\rho$. 
Then, $I_\varepsilon$ limited to $S(a)$ satisfies the $(PS)_c$ condition.
\end{lemma}

\begin{proof}
By Lemma \ref{lem:2.2} $\{u_n\}$ is bounded.  
Let $\{u_n\}\subset S(a)$ be $(PS)_c$ sequence of $I_\varepsilon$ with 
$u_n\rightharpoonup u_\varepsilon$, where $u_\varepsilon\neq 0$ by Lemma \ref{lem:3.2} 
and $c<M_{\rm max}(a)+\rho$. Set $v_n=u_n- u_\varepsilon$. If $v_n\to0$ in $X$, 
then the proof is complete.
If $v_n\not\to0$ in $X$ and $|u_\varepsilon|_{p}=b$, by Lemma \ref{lem:3.3}, 
we obtain
\begin{equation}\label{eqn-(3.14)}
\liminf_{n\to +\infty} |v_n|_p^{p}\ge \beta_1
\end{equation}
for some $\beta_1>0$ which is independent of $\varepsilon\in(0,\varepsilon_0)$.

Let $|v_n|_p=d_n\to d\ge{\beta_{1}}^{1/p}$, by Br\'ezis-Lieb lemma \cite{ref31},  
we have $a^{p}=b^{p}+d^{p}$. From $d_n\in(0,a)$ for $n$ large enough, we have
$$
c+o_n(1)=I_\varepsilon(u_n)=I_\varepsilon(v_n)+I_\varepsilon(u_\varepsilon)+o_n(1)
\ge M_\infty(d_n)+M_{\rm max}(b)+o_n(1).
$$
By Lemma \ref{lem:2.3} (i) and \eqref{eqn-(2.5)}, letting $n\to+\infty$, we have
$$
M_{\rm max}(a)+\rho>c\ge M_\infty(d)+M_{\rm max}(b)
\ge\frac{d^{p}}{a^{p}}M_\infty(a)+\frac{b^{p}}{a^{p}}M_{\rm max}(a).
$$
Then
$$
\rho\ge\frac{d^{p}}{a^{p}}(M_\infty(a)-M_{\rm max}(a))
\ge\frac{\beta_1}{a^{p}}(M_\infty(a)-M_{\rm max}(a)).
$$
This contradicts $\rho<\frac{\beta_1}{a^{p}}(M_\infty(a)-M_{\rm max}(a))$. 
Hence, $v_n \to 0 $ holds in $X$, i.e.\ $u_n \to u_\varepsilon$ in $X$. 
This implies that $u_\varepsilon \in S(a)$ and
\begin{equation*}
(-\Delta)_p^{s} u_\varepsilon+(-\Delta)_q^{s} u_\varepsilon
=\lambda_\varepsilon |u_\varepsilon|^{p-2}u_\varepsilon+h(\varepsilon x)
f(u_\varepsilon)+g(u_\varepsilon),\quad x\in\mathbb{R}^N. \qedhere
\end{equation*}
\end{proof}

\section{Multiplicity result}

In the following, we will discuss some technical issues. 
Let $\rho_0,~r_0> 0,~e_j$ defined in (A6), satisfy
\begin{itemize}
\item $\overline{B_{\rho_{0}}(e_i)}\cap\overline{B_{\rho_{0}}(e_j)}=\emptyset$
for $i\neq j$ and$i,j\in{\{1,\dots,l\}}$;

\item $\cup_{i=1}^{l}B_{\rho_{0}}(e_i)\subset B_{r_{0}}(0)$;

\item $K_{\frac{\rho_0}{2}}=\bigcup_{i=1}^{l}\overline {B_{\frac{\rho_{0}}{2}}(e_i)}$.
\end{itemize}
We set $\kappa:\mathbb{R}^N\to\mathbb{R}^N$ with
$$
\kappa(x)= \begin{cases}
x, &\text{if }|x|\le r_0,\\
r_0\frac{x}{|x|}, &\text{if }|x|>r_0.
\end{cases}
$$
Now, we consider the function $G_\varepsilon : X\backslash\{0\} \to \mathbb{R}^N$ 
defined by
$$
G_\varepsilon(u):=\frac{\int_{\mathbb{R}^N}\kappa(\varepsilon x)|u|^{p}
\mathrm{d}x}{\int_{\mathbb{R}^N}|u|^{p}\mathrm{d}x}.
$$
Then, by the next two theorems, we can obtain the existence of the (PS) sequence
 $I_\varepsilon$ restricted to $S(a)$.

\begin{lemma}\label{lem:4.1}
Reducing $\varepsilon_0$ if necessary, there exists a positive constant 
$\delta_0 < \rho$ such that
$$
G_\varepsilon(u) \in K_{\frac{\rho_0}{2}},\quad\forall \varepsilon\in(0,\varepsilon_0),
$$ 
where $u\in S(a)$ and $I_\varepsilon(u)\le M_{\rm max}(a)+\delta_0$.
\end{lemma}

\begin{proof}
We assume that the conclusion is false, and hence there exist 
$\delta_n \to 0, u_n \in S(a)$ and $\varepsilon_n \to 0$, such that
$$
I_{\varepsilon_n }(u_n)\le M_{\rm max}(a)+\delta_n
$$
and
$G_{\varepsilon_n}(u_n)\not\in K_{\frac{\rho_0}{2}}$.
Firstly, we have
$$
M_{\rm max}(a)\le I_{\rm max}(u_n)\le I_{\varepsilon_n }(u_n)
\le M_{\rm max}(a)+\delta_n,
$$
then
$$
I_{\rm max}(u_n)\to M_{\rm max}(a),\quad \text{as } n\to\infty.
$$
We will analyze the following two cases by means of Lemma \ref{pro:2.1}.

(i) $u_n\to u$ in $X$, where $u\in S(a)$. According to Lebesgue convergence 
theorem, we can deduce that
$$
G_{\varepsilon_n}(u_n)=\frac{\int_{\mathbb{R}^N}\kappa(\varepsilon_n x)|u_n|^{p}
\mathrm{d}x}{\int_{\mathbb{R}^N}|u_n|^{p}\mathrm{d}x}
\to\frac{\int_{\mathbb{R}^N}\kappa(0)|u|^{p}\mathrm{d}x}{\int_{\mathbb{R}^N}
|u|^{p}\mathrm{d}x}=0\in K_{\frac{\rho_0}{2}},
$$
which contradicts to $G_{\varepsilon_n}(u_n)\not\in K_{\frac{\rho_0}{2}}$ 
for $n$ large.

(ii) There exists a sequence $v_n(\cdot)=u_n(\cdot+y_n)$ with $|y_n|\to+\infty$ and 
$\{y_n\}\subset \mathbb{R}^N$ which converges to $v \in S(a)$ in $X$.
Then, we can also study the following two cases:

When $|\varepsilon_n y_n|\to+\infty$, We can infer that
\begin{align*}
 I_{\varepsilon_n}(u_n)
&= \frac{1}{p}\iint_{\mathbb{R}^{2N}}\frac{|v_n(x)-v_n(y)|^{p}}{|x-y|^{N+sp}}
 \mathrm{d}x\mathrm{d}y+\frac{1}{q}\iint_{\mathbb{R}^{2N}}\frac{|v_n(x)
 -v_n(y)|^{q}}{|x-y|^{N+sq}}\mathrm{d}x\mathrm{d}y\\
&\quad -\int_{\mathbb{R}^N}h(\varepsilon_n x+\varepsilon_n y_n)F(v_n)\mathrm{d}x
  -\int_{\mathbb{R}^N}G(v_n)\mathrm{d}x\\
&\to I_\infty(v).
 \end{align*}
Since $I_{\varepsilon_n}(u_n)\le M_{\rm max}(a)+\delta_n$, it holds
$$
M_{\rm max}(a)\ge I_\infty(v)\ge M_\infty(a),
$$
which contradicts to \eqref{eqn-(3.1)}.

When $\varepsilon_n y_n\to y$ for some $y\in\mathbb{R}^N$, we have
\begin{align*}
I_{\varepsilon_n}(u_n)
&= \frac{1}{p}\iint_{\mathbb{R}^{2N}}\frac{|v_n(x)-v_n(y)|^{p}}{|x-y|^{N+sp}}
 \mathrm{d}x\mathrm{d}y+\frac{1}{q}\iint_{\mathbb{R}^{2N}}
 \frac{|v_n(x)-v_n(y)|^{q}}{|x-y|^{N+sq}}\mathrm{d}x\mathrm{d}y\\
&\quad -\int_{\mathbb{R}^N}h(\varepsilon_n x+\varepsilon_n y_n)F(v_n)\mathrm{d}x
 -\int_{\mathbb{R}^N}G(v_n)\mathrm{d}x\\
&\to I_{h(y)}(v),
\end{align*}
then we have
\begin{equation}\label{eqn-(4.1)}
M_{h(y)}(a)\le M_{\rm max}(a).
\end{equation}
Since $h(y) < h_{\rm max}$, Lemma \ref{cor:2.1} implies that $M_{h(y)}
(a) > M_{\rm max}(a)$,
 which contradicts \eqref{eqn-(4.1)}.
Therefore $h(y) = h_{\rm max}$ holds, i.e.\ $y = e_i$
for some $i = 1,\dots,l$. Then we have
\begin{align*}
G_{\varepsilon_n}(u_n)
&=\frac{\int_{\mathbb{R}^N}\kappa(\varepsilon_n x)|u_n|^{p}\mathrm{d}x}
 {\int_{\mathbb{R}^N}|u_n|^{p}\mathrm{d}x}
 =\frac{\int_{\mathbb{R}^N}\kappa(\varepsilon_n x+\varepsilon_n y_n)|v_n|^{p}
 \mathrm{d}x}{\int_{\mathbb{R}^N}|v_n|^{p}\mathrm{d}x}\\
&\to\frac{\int_{\mathbb{R}^N}\kappa(y)|v|^{p}\mathrm{d}x}
{\int_{\mathbb{R}^N}|v|^{p}\mathrm{d}x}=e_i\in K_{\frac{\rho_0}{2}},
\end{align*}
which contradicts to $G_{\varepsilon_n}(u_n)\not\in K_{\frac{\rho_0}{2}}$ for $n$ 
large.
\end{proof}

Next, we introduce some symbols:
\begin{itemize}
\item $\theta_\varepsilon^{i}:=\{u\in S(a):|G_{\varepsilon}(u)-e_i|\le\rho_0\}$,
\item $\partial\theta_\varepsilon^{i}:=\{u\in S(a):|G_{\varepsilon}(u)-e_i|=\rho_0\}$,
\item $\eta_\varepsilon^{i}:=\inf_{u\in\theta_\varepsilon^{i}}  I_\varepsilon(u)$,
\item $\widetilde\eta_\varepsilon^{i}:
=\inf_{u\in\partial\theta_\varepsilon^{i}}  I_\varepsilon(u)$.
\end{itemize}

\begin{lemma}\label{lem:4.2}
Let $0<\delta_0<\rho<\min\{{\frac{1}{p},\frac{\beta}{a^{p}}}\}
(M_\infty(a)-M_{\rm max}(a))$.
Then
$$
\eta_\varepsilon^{i}<M_{\rm max}(a)+\rho\quad \text{and} \quad 
\eta_\varepsilon^{i}<\widetilde\eta_\varepsilon^{i},\quad 
\forall\varepsilon\in(0,\varepsilon_0).
$$
\end{lemma}

\begin{proof}According to Lemmas \ref{pro:2.1} and \ref{lem:2.5}, we set
$$
M_{\rm max}(a)=I_{\rm max}(u),\quad I_{\rm max}'(u)=0,
$$
where $u\in S(a)$. Let $u_\varepsilon^{i}:\mathbb{R}^N\to\mathbb{R}$ be 
$u_\varepsilon^{i}=u(x-e_i/\varepsilon)$ for $1\le i\le l$.
By direct calculation we have
\begin{align*}
I_{\varepsilon}(u_\varepsilon^{i}(x))
&= \frac{1}{p}\iint_{\mathbb{R}^{2N}}\frac{|u(x)-u(y)|^{p}}{|x-y|^{N+sp}}\mathrm{d}x\mathrm{d}y+\frac{1}{q}\iint_{\mathbb{R}^{2N}}\frac{|u(x)-u(y)|^{q}}{|x-y|^{N+sq}}\mathrm{d}x\mathrm{d}y\\
&\quad -\int_{\mathbb{R}^N}h(\varepsilon x+e_i)F(u)\mathrm{d}x-\int_{\mathbb{R}^N}G(u)\mathrm{d}x,
\end{align*}
which implies that
\begin{equation}\label{eqn-(4.2)}
\limsup_{\varepsilon\to0}I_{\varepsilon}(u_\varepsilon^{i}(x))
\le I_{\rm max}(u)=M_{\rm max}(a).
\end{equation}
If $\varepsilon\to0^{+}$, then
$$
G_{\varepsilon}(u_\varepsilon^{i})
=\frac{\int_{\mathbb{R}^N}\kappa(\varepsilon x)|u_\varepsilon^{i}|^{p}\mathrm{d}x}
{\int_{\mathbb{R}^N}|u_\varepsilon^{i}|^{p}\mathrm{d}x}
=\frac{\int_{\mathbb{R}^N}\kappa(\varepsilon x+e_i)|u|^{p}\mathrm{d}x}
{\int_{\mathbb{R}^N}|u|^{p}\mathrm{d}x}\to e_i.
$$
 We can deduce that $u_\varepsilon^{i}\in\theta_\varepsilon^{i}$ when 
$\varepsilon$ is small enough. Moreover, according to \eqref{eqn-(4.2)},
$$
I_{\varepsilon}(u_\varepsilon^{i}(x))
\le M_{\rm max}(a)+\frac{\delta_0}{4},\quad\forall\varepsilon\in(0,\varepsilon_0). 
$$
Hence, reduce $\varepsilon_0$ if necessary,
$$
\eta_\varepsilon^{i}\le M_{\rm max}(a)+\frac{\delta_0}{4}, 
\quad\forall\varepsilon\in(0,\varepsilon_0).
$$
Then
$$
\eta_\varepsilon^{i}\le M_{\rm max}(a)+\rho,\quad
\forall\varepsilon\in(0,\varepsilon_0).
$$

If there exists $u\in\partial\theta_\varepsilon^{i}$, i.e.
$$
u\in S(a)\quad \text{and} \quad |G_{\varepsilon}(u)-e_i|=\rho_0>\frac{\rho_0}{2},
$$
then $G_{\varepsilon}(u)\not\in K_{\frac{\rho_0}{2}}$.
 In conjunction with Lemma \ref{lem:4.1}, we obtain
$$
I_\varepsilon(u)>M_{\rm max}(a)+\frac{\delta_0}{2},\quad
\forall u\in\partial\theta_\varepsilon^{i},~\forall \varepsilon\in(0,\varepsilon_0),
$$
and so,
$$
\widetilde\eta_\varepsilon^{i}>M_{\rm max}(a)+\frac{\delta_0}{2},\quad
\forall\varepsilon\in(0,\varepsilon_0).
$$
From this, it can be seen that
$\eta_\varepsilon^{i}<\widetilde\eta_\varepsilon^{i}$ for all 
$\varepsilon\in(0,\varepsilon_0)$.
\end{proof}

\section{Proof of Theorem \ref{theorem:1.1}}

According to Ekeland's variational principle, we  there exists a sequence 
$\{u_n^{i}\}\subset \theta_\varepsilon^{i}$ $(\subset S(a))$ such that
$$
I_{\varepsilon}(u_n^{i})\to\eta_\varepsilon^{i}
$$
and
$$
I_\varepsilon(v)-I_{\varepsilon}(u_n^{i})\ge-\frac{1}{n}\|v-u_n^{i}\|,\quad
\forall v\in\theta_\varepsilon^{i}\quad \text{with} \quad v\neq u_n^{i}
$$
for each $i\in\{1,\dots,l\}$.

Then, getting $u_n^{i}\in\theta_\varepsilon^{i}\backslash\partial\theta_\varepsilon^{i}$
 for sufficiently large $n$ by Lemma \ref{lem:4.2}.
Given $v\in T_{u_n^{i}}S(a)=\{\omega\in X:\int_{\mathbb{R}^N}
|u_n^{i}|^{p-2}u_n^{i}\omega \mathrm{d}x=0\}$,
We can define the path $\sigma:(-\xi_1,\xi_1)\to S(a)$
with
$$
\sigma(t)=a\frac{(u_n^{i}+tv)}{|u_n^{i}+tv|_p},
$$
where $\xi_1>0$. It is clear that $\sigma\in C^{1}((-\xi_1,\xi_1),S(a))$,
and
$$
\sigma(t)\in\theta_\varepsilon^{i}\backslash\partial\theta_\varepsilon^{i},\quad 
\forall t\in(-\xi_1,\xi_1),\quad \sigma(0)=u_n^{i}\quad \text{and} \quad \sigma'(0)=v.
$$
Then we have
$$
I_\varepsilon(\sigma(t))-I_\varepsilon(u_n^{i})\ge-\frac{1}{n}\|\sigma(t)-u_n^{i}\|
$$
for $t\in(-\xi_1,\xi_1)$, this means that
\begin{align*}
\frac{I_\varepsilon(\sigma(t))-I_\varepsilon(\sigma(0))}{t}
&=\frac{I_\varepsilon(\sigma(t))-I_\varepsilon(u_n^{i})}{t}\\
&\ge-\frac{1}{n}\|\frac{\sigma(t)-u_n^{i}}{t}\|\\
&=-\frac{1}{n}\|\frac{\sigma(t)-\sigma(0)}{t}\|,\quad\forall t\in(0,\xi_1).
\end{align*}
Taking the limit of $t\to0^{+}$, we obtain
$$
\langle I_{\varepsilon}' (u_n^{i}),v \rangle\ge-\frac{1}{n}\|v\|.
$$
Then, instead of $v$, we can use $-v$ to derive
$$
\sup\{|\langle I_{\varepsilon}' (u_n^{i}),v \rangle|:\|v\|\le1\}\le\frac{1}{n},
$$
which implies that
$$
I_{\varepsilon}(u_n^{i})\to\eta_\varepsilon^{i}\quad \text{and}\quad
\|I_\varepsilon|_{S(a)}'(u_n)\|\to 0\quad \text{as} \quad n\to\infty, 
$$
which means $\{u_n^{i}\}\subset S(a)$ is a $(PS)_{\eta_\varepsilon^{i}}$
sequence of $I_\varepsilon$.
Combining Lemma \ref{lem:3.4} and $\eta_\varepsilon^{i}<M_{\rm max}(a)+\rho$, we can 
deduce that there exists $u^{i}$ such that $u_n^{i}\to u^{i}$ in $X$,
$$
u^{i}\in\theta_\varepsilon^{i},~ I_\varepsilon(u^{i})=\eta_\varepsilon^{i}\quad 
\text{and}\quad I_\varepsilon|_{S(a)}'(u^{i})=0.
$$
According to our assumptions, we have
$$
G_\varepsilon(u^{i})\in\overline{B_{\rho_{0}}(e_i)},\quad 
G_\varepsilon(u^{j})\in\overline{B_{\rho_{0}}(e_j)}
$$
and
$$
\overline{B_{\rho_{0}}(e_i)}\cap\overline{B_{\rho_{0}}(e_j)}=\emptyset \quad 
\text{for }  i\neq j  \text{ and }i,j\in{\{1,\dots,l\}},
$$
which means $u^{i}\neq u^{j}$ for $i \neq  j$ while $1 \le i, j \le l$.
Thus, for any $\varepsilon\in(0,\varepsilon_0)$, $I_\varepsilon$ has at least 
$l$ nontrivial critical points $(u^{i},\lambda_i)$, i.e.
$$
(-\Delta)_p^{s} u^{i}+(-\Delta)_q^{s} u^{i}=\lambda_i |u^{i}|^{p-2}u^{i}
+h(\varepsilon x)f(u^{i})+g(u^{i}),\quad \forall i\in\{1,2,\dots,l\},
$$
which means
\begin{align*}
&\iint_{\mathbb{R}^{2N}}\frac{|u^{i}(x)-u^{i}(y)|^{p}}{|x-y|^{N+sp}}\mathrm{d}x\mathrm{d}y+\iint_{\mathbb{R}^{2N}}\frac{|u^{i}(x)-u^{i}(y)|^{q}}{|x-y|^{N+sq}}\mathrm{d}x\mathrm{d}y-\int_{\mathbb{R}^N}\lambda_{i} |u^{i}|^{p}\mathrm{d}x\\
&-\int_{\mathbb{R}^N}h(\varepsilon x) f(u^{i})u^{i}\mathrm{d}x-\int_{\mathbb{R}^N} g(u^{i})u^{i}\mathrm{d}x=0.
\end{align*}
Combining with $I_\varepsilon(u^{i})<0$, we have
\begin{align*}
0&>I_\varepsilon(u^{i})-\frac{1}{q}(\iint_{\mathbb{R}^{2N}}\frac{|u^{i}(x)
 -u^{i}(y)|^{p}}{|x-y|^{N+sp}}\mathrm{d}x\mathrm{d}y
 +\iint_{\mathbb{R}^{2N}}\frac{|u^{i}(x)-u^{i}(y)|^{q}}{|x-y|^{N+sq}}
 \mathrm{d}x\mathrm{d}y-\int_{\mathbb{R}^N}\lambda_{i} |u^{i}|^{p}\mathrm{d}x\\
&\quad -\int_{\mathbb{R}^N}\mu f(u^{i})u^{i}\mathrm{d}x
 -\int_{\mathbb{R}^N} g(u^{i})u^{i}\mathrm{d}x)\\
&= (\frac{1}{p}-\frac{1}{q})\iint_{\mathbb{R}^{2N}}\frac{|u^{i}(x)-u^{i}(y)|^{p}}
 {|x-y|^{N+sp}}\mathrm{d}x\mathrm{d}y
 +\frac{1}{q}\int_{\mathbb{R}^N}\lambda_i |{u^{i}}|^{p}\mathrm{d}x\\
&\quad +\frac{1}{q}\int_{\mathbb{R}^N}h(\varepsilon x)
 f(u^{i})u^{i}\mathrm{d}x-\int_{\mathbb{R}^N}h(\varepsilon x)  F(u^{i})\mathrm{d}x\\
&\quad +\frac{1}{q}\int_{\mathbb{R}^N} g(u^{i})u^{i}\mathrm{d}x
 -\int_{\mathbb{R}^N} G(u^{i})\mathrm{d}x\\
&\geq \frac{1}{q}\int_{\mathbb{R}^N}\lambda_i |{u^{i}}|^{p}\mathrm{d}x,
\end{align*}
which implies $\lambda_i<0$. This proves the expected result.

\subsection*{Acknowledgements}
This work was supported by the Natural Science Foundation of Jilin Province 
(Grant No. 20210101156JC).


\begin{thebibliography}{10}

\bibitem{ref1}
V. Ambrosio, T. Isernia;
Multiplicity and concentration results for some nonlinear Schr\"{o}dinger equations 
with the fractional $p$-Laplacian,
\emph{Discrete Contin. Dyn. Syst.}, \textbf{38} (2018), 5835-5881.

\bibitem{ref2}
V. Ambrosio, T. Isernia;
On a fractional $p\&q$ Laplacian problem
with Critical Sobolev-Hardy Exponents,
\emph{Mediterr. J. Math.},
\textbf{15} (2018), 1-17.

\bibitem{ref3}
V. Ambrosio;
Nontrivial solutions for a fractional $p$-Laplacian problem via Rabier Theorem,
\emph{Complex Var. Elliptic Equations},
\textbf{62} (2017), 838-847.

\bibitem{ref4}
V. Ambrosio;
Multiple solutions for a fractional $p$-Laplacian equation with sign-changing potential,
\emph{Electron. J. Differ. Equations},
\textbf{2016} (2016), 1-12.

\bibitem{n5} 
V. Ambrosio; 
Fractional p\&q Laplacian Problems in $\mathbb{R}^N$  with Critical Growth. 
\emph{Z. Anal. Anwend.} \textbf{39}(2020), 289–314

\bibitem{ref5}
F. J. Almgren, Jr.,  E. H. Lieb;
Symmetric decreasing rearrangement is sometimes continuous,
\emph{J. Am. Math. Soc.},
\textbf{2} (1989), 683-773.

\bibitem{ref6}
J. G. Azorero, I. P. Alonso;
Multiplicity of solutions
for elliptic problems with critical exponent
or with a nonsymmetric term,
\emph{Trans. Am. Math. Soc.},
\textbf{323} (1991), 20 pp.


\bibitem{ref7}
G. M. Bisci, V. D. R\u{a}dulescu;
Ground state solutions of scalar field fractional
Schr\"{o}dinger equations,
\emph{Calc. Var.},
\textbf{54} (2015), 2985-3008.

\bibitem{ref8}
A. Bahrouni, Q. Guo, H. Hajaiej;
Normalized solutions to the mixed fractional Schr\"{o}dinger equations with potential and general nonlinear term,
\emph{Asymptotic. Anal.},
\textbf{} (2025), 30 pp.

\bibitem{ref9}
N. Biswas, F. Sk;
On generalized eigenvalue problems of fractional $(p, q)$-Laplace operator with two parameters,
\emph{P. Roy. Soc. Edinb. A.},
\textbf{} (2022), 1-46.

\bibitem{n1} J. Chaker, M. Kim, and M. Weidner; Harnack inequality for nonlocal problems with
non-standard growth. \emph{Math. Ann.} \textbf{386} (2023), 533–550.

\bibitem{n2} J. Chen,  Y. Gao;
Multiple normalized solutions for the non-autonomous $p$-Laplace equation, \emph{Communications on Pure and Applied Analysis}, \textbf{23} (2024), 1044-1063.

\bibitem{ref10}
M. Chhetri, P. Girg, E. Hollifield;
Existence of positive solutions for fractional Laplacian equations: 
theory and numerical experiments,
\emph{Electron. J. Differ. Equations},
\textbf{2020} (2020), 1-31.

\bibitem{n3} X. Cui, A. Li, and C. Wei;
 Existence of normalized solutions to a class of non-autonomous
$(p, q)$-Laplacian equations, \emph{Anal. Math. Phys.}, \textbf{15} (2025), 30 pp.


\bibitem{ref11}
S. Dipierro, G. Palatucci, E. Valdinoci;
Existence and symmetry results for a Schr\"{o}dinger type problem involving the 
fractional Laplacian, \emph{Mathematiche}, \textbf{68} (2013), 201-216.

\bibitem{ref12}
Y. Ding, HY. Wang;
Normalized Solutions to Schr\"{o}dinger equations with critical exponent and mixed 
nonlocal nonlinearities,
\emph{ J. Geom. Anal.}, \textbf{215} (2024), 38 pp.


\bibitem{ref13}
I. Ekeland;
On the variational principle,
\emph{J. Math. Anal. Appl.}, \textbf{47} (1974), 324-353.

\bibitem{ref14}
P. Felmer, A. Quaas, J. Tan;
Positive solutions of the nonlinear Schr\"{o}dinger
equation with the fractional Laplacian,
\emph{P. Roy. Soc. Edinb. A.}, \textbf{142} (2012), 1237-1262.

\bibitem{ref15}
G. M. Figueiredo;
Existence of positive solutions for a class of $p\&q$ elliptic problems with 
critical growth on $\mathbb{R}^N$,
\emph{J. Math. Anal. Appl.}, \textbf{378} (2011), 507-518.


\bibitem{n4} 
J. Giacomoni, D. Kumar,  K. Sreenadh;
 Global regularity results for nonhomogeneous growth fractional problems. 
 \emph{J. Geom. Anal.}, \textbf{32}(2022), Paper No. 36.

\bibitem{ref16}
D. Goel, D. Kumar, K. Sreenadh;
Regularity and multiplicity results for fractional
$(p, q)$-Laplacian equations,
\emph{Commun. Contemp. Math.}, \textbf{22} (2020), 36 pp.

\bibitem{ref18}
A. Iannizzotto, R. Livrea;
Four Solutions for fractional $p$-Laplacian
equations with Asymmetric Reactions,
\emph{Mediterr. J. Math.}, \textbf{18} (2021), 32 pp.

\bibitem{ref19}
H. Luo, Z. Zhang;
Normalized solutions to the fractional Schr\"{o}dinger equations with combined
 nonlinearities, \emph{Calc. Var.},
\textbf{59} (2020), 35 pp.

\bibitem{ref20}
Q. Li , W. Zou;
The existence and multiplicity of the
normalized solutions for fractional
Schr\"{o}dinger equations involving Sobolev
critical exponent in the $L^2$-subcritical and
$L^2$-supercritical cases,
\emph{Adv. Nonlinear Anal.}, \textbf{11} (2021), 1531-1551.

\bibitem{ref21}
G. Li, X. Liang;
The existence of nontrivial solutions to nonlinear elliptic equation of 
$p-q$-Laplacian type on $\mathbb{R}^N$,
\emph{Nonlinear Anal.}, \textbf{71} (2009), 2316-2334.

\bibitem{ref22}
E. H. Lieb, M. Loss;
Analysis, Graduate Studies in Mathematics,
\emph{American Mathematical Society}, Providence, RI, 2001.

\bibitem{ref23}
T. V. Nguyen, V. D. R\u{a}dulescu;
Multiple normalized solutions for fractional
elliptic problems,
\emph{Forum Math.}, \textbf{36} (2024), 1225-1248.

\bibitem{ref24}
H. Nguyen,  M. Squassina;
Fractional Caffarelli-Kohn-Nirenberg inequalities,
\emph{J. Funct. Anal.}, \textbf{274} (2018), 2661-2672.

\bibitem{ref25}
S. Peng, A. Xia;
Normalized solutions of supercritical nonliner fractional
Schr\"{o}dinger equation with potential,
\emph{Commun. Pure Appl. Anal.}, \textbf{20} (2021), 3723-3744.

\bibitem{ref26}
P. Pucci, M. Xiang, B. Zhang;
Multiple solutions for nonhomogeneous
Schr\"{o}dinger-Kirchhoff type equations involving
the fractional $p$-Laplacian in $\mathbb{R}^N$,
\emph{Calc. Var.}, \textbf{54} (2015), 2785-2806.

\bibitem{ref27}
E. D. Silva, J. L. A. Oliveira, C. Goulart;
Quasilinear nonlocal elliptic problems with prescribed norm in the $ L^{p}$-subcritical 
and $ L^{p}$-critical growth, {arXiv: 2412.14948}.

\bibitem{ref28}
N. Soave;
Normalized ground states for the NLS equation with combined nonlinearities,
\emph{J. Differ. Equ.}, \textbf{269} (2020), 6941-6987.

\bibitem{ref29}
N. Soave;
Normalized ground states for the NLS equation with combined nonlinearities: 
the Sobolev critical case,
\emph{J. Funct. Anal.}, \textbf{279} (2020), 43 pp.

\bibitem{ref30}
N. V. Thin, P. T. Thuy, T. T. D. Linh;
Existence of solution for the $(p,q)$-fractional
Laplacian equation with nonlocal Choquard
reaction and exponential growth,
\emph{Complex Var. Elliptic Equations}, \textbf{69} (2024), 1949-1972.

\bibitem{ref31}
M. Willem;
Minimax theorems, Progress in nonlinear differential equations and their applications,
\emph{Birkh\"{a}user Boston, Inc.}, Boston, MA, 1996.

\bibitem{ref32}
M. Xiang, B. Zhang, X. Zhang;
A nonhomogeneous fractional $p$-Kirchhoff
type problem involving critical exponent
in $\mathbb{R}^N$,
\emph{Adv. Nonlinear Stud.}, \textbf{17} (2017), 611-640.

\bibitem{ref33}
M. Xiang, Y. Ma, M. Yang;
Normalized homoclinic solutions of discrete nonlocal double phase problems,
\emph{Bull. Math. Sci}, \textbf{14} (2024), 18 pp.

\bibitem{ref34}
X. Zhang, M. Squassina, J. Zhang;
Multiplicity of Normalized Solutions for the Fractional
Schr\"{o}dinger Equation with Potentials,
\emph{Mathematics}, \textbf{12} (2024), 20 pp.

\bibitem{ref35}
X. Zhang, S. Liang;
Multiple normalized solutions to a Choquard equation
involving fractional $p$-Laplacian in $\mathbb{R}^N$,
\emph{Fractal Fract.}, \textbf{8} (2024), 27 pp.

\bibitem{ref36}
X. Zhang, T. V. Nguyen,  S. Liang;
Multiple normalized solutions to critical Choquard equation involving fractional $p$-Laplacian in $\mathbb{R}^N$,
\emph{Anal. Math. Phys.}, \textbf{15} (2025), 42 pp.

\bibitem{ref37}
E. WB. Zongo, P. A. Feulefack;
Bifurcation results and multiple solutions for the fractional 
$(p, q) $-Laplace operators,
{arXiv:2406.15825}.

\end{thebibliography}

\end{document}
